\documentclass[11pt,a4paper]{article}

% ============================================================
% PACOTES
% ============================================================
\usepackage[utf8]{inputenc}
\usepackage[T1]{fontenc}
\usepackage[brazilian]{babel}

\usepackage{amsmath,amssymb,amsfonts,mathtools}
\usepackage{bm}
\usepackage{physics}
\usepackage{siunitx}
\usepackage{booktabs}
\usepackage{geometry}
\usepackage{float}
\usepackage{tikz}
\usepackage{pgfplots}
\usepackage{listings}
\usepackage{xcolor}
\usepackage{hyperref}
\usetikzlibrary{intersections}

\pgfplotsset{compat=1.18}
\geometry{margin=2.5cm}

\hypersetup{
    colorlinks=true,
    linkcolor=blue,
    citecolor=blue,
    urlcolor=blue
}

% ============================================================
% COMANDOS
% ============================================================
\newcommand{\pdvT}[2]{\left(\frac{\partial #1}{\partial #2}\right)_T}
\newcommand{\pdvV}[2]{\left(\frac{\partial #1}{\partial #2}\right)_v}
\newcommand{\pdvS}[2]{\left(\frac{\partial #1}{\partial #2}\right)_s}
\newcommand{\pd}[2]{\frac{\partial #1}{\partial #2}}
\newcommand{\pdd}[2]{\frac{\partial^2 #1}{\partial #2^2}}
\newcommand{\e}[1]{\mathrm{e}^{#1}}
\renewcommand{\dd}{\mathrm{d}}
\newcommand{\avg}[1]{\left\langle #1\right\rangle}

% ============================================================
% DOCUMENTO
% ============================================================
\begin{document}
%
\begin{center}
    {\Large \textbf{Segunda série de exercícios}}\\[0.6cm]
    {\large Tópicos de Mecânica Estatística -- Transições de Fases -- 2026}\\[0.2cm]
    {\large Modelo de Ising e teoria de Landau -- 25/8/2026}\\[0.4cm]
    {\large Data de Entrega -- 01/9/2026}\\[0.2cm]
    {\Large \textbf{Autor:} Thiago Siqueira Domingues (Nº USP: 8942194)}
\end{center}
%
\vspace{0.3cm}

\noindent\textit{Observação sobre as figuras geradas:} todas as figuras foram geradas por mim utilizando a linguagem de programação Python. Os scripts podem ser facilmente encontrados na página do Github~\cite{domingues2026statmech}. 
%
\vspace{0.3cm}
% ============================================================
\section*{Problema (1) -- Modelo de Ising unidimensional}
%
O hamiltoniano de spin do modelo de Ising ferromagnético com interações de primeiros vizinhos em uma rede unidimensional de $N$ sítios, com condições periódicas de contorno ($\sigma_{N+1}\equiv\sigma_1$), é
%
\begin{equation}
    \mathcal H = -J\sum_{i=1}^N \sigma_i\sigma_{i+1} - H\sum_{i=1}^N\sigma_i,
    \qquad J>0,\quad \sigma_i=\pm1.
    \label{eq:ising_H}
\end{equation}
%
\subsection*{(a) Matriz de transferência e correlação de pares em campo nulo}
%
Seguindo Ref.~\cite{salinas}, primeiro temos que calcular a função canônica de partição:
%
\begin{equation}
Z_N = \sum_{\{\sigma_i\}} \exp\left( -\beta \mathcal H \right)
= \sum_{\{\sigma_i\}} \exp\left( K \sum_{i=1}^N \sigma_i\sigma_{i+1} + L \sum_{i=1}^N\sigma_i \right),
\end{equation}
%
onde definimos os parâmetros adimensionais
%
\begin{equation}
K \equiv \beta J, \qquad L \equiv \beta H.
\end{equation}
%
Reescrevendo o segundo termo de maneira mais simétrica:
%
\begin{equation}
L \sum_{i=1}^N \sigma_i
= \sum_{i=1}^N \frac{L}{2}(\sigma_i + \sigma_{i+1}),
\end{equation}
%
que é válido devido às condições periódicas de contorno (\(\sigma_{N+1}=\sigma_1\)). Assim, podemos escrever a função de partição na forma, 
%
\begin{equation}
Z_N = \sum_{\sigma_1} \sum_{\sigma_2} \cdots \sum_{\sigma_N}
\prod_{i=1}^N \exp\left( K \sigma_i\sigma_{i+1} + \frac{L}{2}(\sigma_i + \sigma_{i+1}) \right).
\end{equation}
%
Definimos a \textbf{matriz de transferência} \(T\) com elementos
%
\begin{equation}
T_{\sigma_i, \sigma_{i+1}}
= \exp\left( K \sigma_i\sigma_{i+1} + \frac{L}{2}(\sigma_i + \sigma_{i+1}) \right).
\end{equation}
%
Na base \(\{\sigma = +1, -1\}\), a matriz é escrita explicitamente como
%
\begin{equation}
T =
\begin{pmatrix}
T_{++} & T_{+-} \\
T_{-+} & T_{--}
\end{pmatrix}
=
\begin{pmatrix}
e^{K+L} & e^{-K} \\
e^{-K} & e^{K-L}
\end{pmatrix}.
\label{eq:transfer_matrix}
\end{equation}
%
Portanto, a função de partição é o traço da \(N\)-ésima potência da matriz de transferência:
%
\begin{equation}
Z_N = \sum_{\sigma_1} \sum_{\sigma_2} \cdots \sum_{\sigma_N} T_{\sigma_1,\sigma_2} T_{\sigma_2,\sigma_3} \cdots T_{\sigma_N,\sigma_1}
= \sum_{\sigma_1} \left( T^N \right)_{\sigma_1, \sigma_1}
= \operatorname{Tr}\left( T^N \right).
\end{equation}
%
Pois temos N matrizes idênticas e utilizamos o produto matricial para escrever a segunda igualdade.
Logo,
%
\begin{equation}
Z_N = \operatorname{Tr}\left( T^N \right).
\end{equation}
%
Como a matriz de transferência \eqref{eq:transfer_matrix} é simétrica, por construção, isso implica que ela é diagonalizável por meio de uma transformação unitária. Assim, existe uma matriz unitária $U$ tal que $U^\dagger U = I$ e
%
\begin{equation}
U^\dagger T U = D,
\end{equation}
%
onde $D$ é uma matriz diagonal contendo os autovalores de $T$. Equivalentemente, podemos escrever $T = U D U^\dagger$. Portanto, a $N$-ésima potência da matriz de transferência é dada por
%
\begin{equation}
T^N = (U D U^\dagger)^N = U D^N U^\dagger.
\end{equation} 
%
A função de partição é o traço de $T^N$:
%
\begin{equation}
Z_N = \operatorname{Tr}(T^N)
= \operatorname{Tr}(U D^N U^\dagger).
\end{equation}
%
Usando a propriedade cíclica do traço ($\operatorname{Tr}(ABC) = \operatorname{Tr}(BCA)$), obtemos
%
\begin{equation}
Z_N = \operatorname{Tr}(D^N U^\dagger U) = \operatorname{Tr}(D^N).
\end{equation}
%
Como $D = \operatorname{diag}(\lambda_1, \lambda_2)$, temos $D^N = \operatorname{diag}(\lambda_1^N, \lambda_2^N)$. Logo,
%
\begin{equation}
Z_N = \lambda_1^N + \lambda_2^N.
\label{eq:Z_autovalores_spin_half}
\end{equation}
%
Agora, calculamos os autovalores $\lambda$ pelas raízes da equação secular,
%
\begin{equation}
\det(T - \lambda I) = 0,
\end{equation}
%
temos
%
\begin{equation}
\det(T - \lambda I) =
\begin{vmatrix}
e^{K+L} - \lambda & e^{-K} \\
e^{-K} & e^{K-L} - \lambda
\end{vmatrix}
= 0.
\end{equation}
%
Calculando o determinante:
%
\begin{equation}
(e^{K+L} - \lambda)(e^{K-L} - \lambda) - e^{-2K} = 0.
\end{equation}
%
Expandindo:
%
\begin{equation}
\lambda^2 - \lambda (e^{K+L} + e^{K-L}) + e^{2K} - e^{-2K} = 0.
\end{equation}
%
Reescrevendo os termos:
%
\begin{equation}
\lambda^2 - 2e^K \lambda\cosh L \,  + 2\sinh(2K) = 0,
\end{equation}
%
onde usamos $e^{K+L} + e^{K-L} = 2e^K \cosh L$ e $e^{2K} - e^{-2K} = 2\sinh(2K)$.
%
Resolvendo a equação quadrática:
%
\begin{equation}
\lambda_\pm = \frac{2e^K \cosh L \pm \sqrt{4e^{2K}\cosh^2 L - 8\sinh(2K)}}{2}.
\end{equation}
%
Simplificando:
%
\begin{equation}
\lambda_\pm = e^K \cosh L \pm \sqrt{e^{2K}\cosh^2 L - 2\sinh(2K)}.
\end{equation}
%
Simplificando o discriminante:
%
\begin{equation}
e^{2K}\cosh^2 L - 2\sinh(2K)
= e^{2K}\cosh^2 L - (e^{2K} - e^{-2K})
= e^{2K}(\cosh^2 L - 1) + e^{-2K}
= e^{2K}\sinh^2 L + e^{-2K}.
\end{equation}
%
Portanto, os autovalores são
%
\begin{equation}
\lambda_\pm
=
e^K \cosh L
\pm
\sqrt{ e^{2K}\sinh^2 L + e^{-2K}}.
\label{eq:autovalores_finais}
\end{equation}
%
Podemos mostrar que os autovalores são sempre positivos e que $\lambda_+ > \lambda_-$ (exceto em $T=H=0$). Observamos que $\lambda_+ > \lambda_-$ para todos os valores de $K > 0$ (temperatura finita) e $L$ real, exceto em um ponto particular. De fato, como o discriminante é sempre positivo,
%
\begin{equation}
e^{2K}\sinh^2 L + e^{-2K} > 0,
\end{equation}
%
temos $\lambda_+ > \lambda_-$ para qualquer $K > 0$. A única exceção ocorre no limite $K \to \infty$ (ou seja, $T \to 0$) com $H = 0$, onde os autovalores se tornam degenerados:
%
\begin{equation}
\lambda_+ = e^K + e^{-K} \longrightarrow e^K, \qquad
\lambda_- = e^K - e^{-K} \longrightarrow e^K,
\end{equation}
%
de modo que $\lambda_+ = \lambda_- = e^K$ no limite $K \to \infty$ (temperatura zero). Para o caso particular de campo nulo ($H = 0 \Rightarrow L = 0$), os autovalores simplificam-se para
%
\begin{equation}
\lambda_+ = e^K + e^{-K} = 2\cosh K, \qquad
\lambda_- = e^K - e^{-K} = 2\sinh K.
\label{eq:autovalores_h0}
\end{equation}
%
Neste caso, temos $\lambda_+ > \lambda_-$ para todo $K > 0$, com $\lambda_+ \to \lambda_-$ apenas no limite $K \to \infty$ (zero absoluto). 
No limite termodinâmico ($N \to \infty$), a função de partição é dominada pelo maior autovalor $\lambda_+$:
%
\begin{equation}
Z_N = \lambda_+^N + \lambda_-^N
= \lambda_+^N \left[ 1 + \left( \frac{\lambda_-}{\lambda_+} \right)^N \right].
\end{equation}
%
Como $\lambda_-/\lambda_+ < 1$, o termo entre colchetes tende a 1 quando $N \to \infty$. Portanto,
%
\begin{equation}
Z_N \simeq \lambda_+^N.
\end{equation}
%
A energia livre de Gibbs por sítio é definida como
%
\begin{equation}
g(T,H) = \lim_{N \to \infty} \left[ -\frac{1}{\beta N} \ln Z_N \right].
\end{equation}
%
Substituindo $Z_N$:
%
\begin{equation}
g(T,H) = -\frac{1}{\beta} \ln \lambda_+
= -\frac{1}{\beta} \ln \left[ e^K \cosh L + \sqrt{ e^{2K}\sinh^2 L + e^{-2K} } \right].
\end{equation}
%
Em termos das variáveis originais $J$, $H$ e $T$:
%
\begin{equation}
g(T,H) = -k_B T \ln \left[ e^{\beta J} \cosh(\beta H) + \sqrt{ e^{2\beta J}\sinh^2(\beta H) + e^{-2\beta J} } \right].
\label{eq:gibbs_free_energy}
\end{equation}
%
Esta energia livre é uma função analítica de $T$ e $H$ para todo $T > 0$, o que implica que não há transição de fase em temperatura finita no modelo de Ising 1D. 
%
A função de correlação de pares é definida como o valor esperado
%
\begin{equation}
\langle \sigma_k \sigma_l \rangle
= \frac{1}{Z_N} \sum_{\{\sigma_i\}} \sigma_k \sigma_l \, e^{-\beta \mathcal H}.
\end{equation}
% 
Agora, para calcular o valor esperado de um produto de spins em dois sítios distintos, digamos \(\sigma_k\) e \(\sigma_l\) (com \(l > k\)), precisamos inserir os fatores \(\sigma_k\) e \(\sigma_l\) dentro da soma sobre configurações. Na linguagem matricial, o fator \(\sigma_i\) é representado pela matriz de spin \(\Sigma\), cujos elementos são,
%
\begin{equation}
\Sigma_{\sigma,\sigma'} = \sigma \, \delta_{\sigma,\sigma'}.
\end{equation}
%
Na base \(\{\sigma = +1, -1\}\), temos \(\Sigma = \operatorname{diag}(1, -1)\), é a matriz de Pauli $\sigma_z$. Ao percorrer a cadeia de sítios, a inserção de \(\sigma_k\) e \(\sigma_l\) no produto de matrizes equivale a colocar uma matriz \(\Sigma\) na posição \(k\) e outra matriz \(\Sigma\) na posição \(l\) dentro do traço. Como a ordem do produto é \(T_{\sigma_1,\sigma_2} T_{\sigma_2,\sigma_3} \cdots T_{\sigma_N,\sigma_1}\), a distância entre \(k\) e \(l\) ao longo da cadeia (no sentido direto) é \(l - k\), e a distância restante para fechar o contorno periódico é \(N - (l - k)\). Portanto,
%
\begin{equation}
\langle \sigma_k \sigma_l \rangle
= \frac{1}{Z_N} \sum_{\{\sigma_i\}} \sigma_k \sigma_l \prod_{i=1}^N T_{\sigma_i,\sigma_{i+1}}
= \frac{ \operatorname{Tr}\left( \Sigma \, T^{l-k} \, \Sigma \, T^{N-l+k} \right) }{ \operatorname{Tr}\left( T^N \right) }.
\label{eq:traco_geral}
\end{equation}
%
Partimos da expressão \eqref{eq:traco_geral} para o caso de campo nulo ($H=0 \Rightarrow L=0$). Para $H=0$, a matriz de transferência \eqref{eq:transfer_matrix} reduz-se a
%
\begin{equation}
T =
\begin{pmatrix}
e^{K} & e^{-K} \\
e^{-K} & e^{K}
\end{pmatrix}.
\end{equation}
%
Seus autovalores, já calculados em \eqref{eq:autovalores_h0}, são
%
\begin{equation}
\lambda_+ = 2\cosh K, \qquad \lambda_- = 2\sinh K,
\end{equation}
%
com autovetores ortonormais
%
\begin{equation}
v_+ = \frac{1}{\sqrt{2}}
\begin{pmatrix}
1 \\
1
\end{pmatrix}, \qquad
v_- = \frac{1}{\sqrt{2}}
\begin{pmatrix}
1 \\
-1
\end{pmatrix}.
\end{equation}
%
Assim, a matriz unitária (ortogonal) que diagonaliza $T$ é
%
\begin{equation}
U = 
\begin{pmatrix}
v_+ & v_-
\end{pmatrix}
= \frac{1}{\sqrt{2}}
\begin{pmatrix}
1 & 1 \\
1 & -1
\end{pmatrix}.
\end{equation}
%
Temos $T = U D U^\dagger$, com $D = \operatorname{diag}(\lambda_+, \lambda_-)$.
%
Agora, definimos a matriz
%
\begin{equation}
A \equiv U^\dagger \Sigma U,
\end{equation}
%
onde $\Sigma = \operatorname{diag}(1, -1)$. Calculando $A$ explicitamente:
%
\begin{equation}
A = \frac{1}{2}
\begin{pmatrix}
1 & 1 \\
1 & -1
\end{pmatrix}
\begin{pmatrix}
1 & 0 \\
0 & -1
\end{pmatrix}
\begin{pmatrix}
1 & 1 \\
1 & -1
\end{pmatrix}.
\end{equation}
%
Primeiro, $\Sigma U = \frac{1}{\sqrt{2}}
%
\begin{pmatrix}
1 & 1 \\
-1 & 1
\end{pmatrix}$. Portanto,
%
\begin{equation}
A = \frac{1}{2}
\begin{pmatrix}
1 & 1 \\
1 & -1
\end{pmatrix}
\begin{pmatrix}
1 & 1 \\
-1 & 1
\end{pmatrix}
= \frac{1}{2}
\begin{pmatrix}
0 & 2 \\
2 & 0
\end{pmatrix}
= \begin{pmatrix}
0 & 1 \\
1 & 0
\end{pmatrix}.
\label{eq:A_matrix}
\end{equation}
%
Logo, os elementos de $A$ são
%
\begin{equation}
A_{++} = 0, \qquad A_{--} = 0, \qquad A_{+-} = A_{-+} = 1.
\end{equation}
%
Substituindo $T^{l-k} = U D^{l-k} U^\dagger$ e $T^{N-l+k} = U D^{N-l+k} U^\dagger$ no traço de \eqref{eq:traco_geral}, obtemos
%
\begin{align}
\operatorname{Tr}\left( \Sigma T^{l-k} \Sigma T^{N-l+k} \right)
&= \operatorname{Tr}\left( \Sigma U D^{l-k} U^\dagger \Sigma U D^{N-l+k} U^\dagger \right) \nonumber \\
&= \operatorname{Tr}\left( U^\dagger \Sigma U D^{l-k} U^\dagger \Sigma U D^{N-l+k} \right) \nonumber \\
&= \operatorname{Tr}\left( A D^{l-k} A D^{N-l+k} \right),
\end{align}
%
onde na segunda igualdade usamos a propriedade cíclica do traço.
%
Agora, expandimos o traço explicitamente:
%
\begin{equation}
\operatorname{Tr}\left( A D^{l-k} A D^{N-l+k} \right)
= \sum_{i,j} A_{ij} A_{ji} \, \lambda_i^{l-k} \, \lambda_j^{N-l+k},
\end{equation}
%
com $i,j \in \{+, -\}$.
%
Como $A_{++} = A_{--} = 0$, os únicos termos não nulos são aqueles com $i \neq j$:
%
\begin{itemize}
\item Para $i = +$ e $j = -$:
\begin{equation}
A_{+-} A_{-+} \, \lambda_+^{l-k} \lambda_-^{N-l+k}
= 1 \cdot 1 \cdot \lambda_+^{l-k} \lambda_-^{N-l+k}.
\end{equation}
%
\item Para $i = -$ e $j = +$:
\begin{equation}
A_{-+} A_{+-} \, \lambda_-^{l-k} \lambda_+^{N-l+k}
= 1 \cdot 1 \cdot \lambda_-^{l-k} \lambda_+^{N-l+k}.
\end{equation}
\end{itemize}
%
Somando as duas contribuições:
%
\begin{equation}
\operatorname{Tr}\left( \Sigma T^{l-k} \Sigma T^{N-l+k} \right)
= \lambda_+^{l-k} \lambda_-^{N-l+k} + \lambda_-^{l-k} \lambda_+^{N-l+k}.
\label{eq:trace_corr_calculado}
\end{equation}
%
Reorganizando a ordem dos fatores no primeiro termo (apenas por conveniência de notação, pois a multiplicação é comutativa), podemos escrever
%
\begin{equation}
\operatorname{Tr}\left( \Sigma T^{l-k} \Sigma T^{N-l+k} \right)
= \lambda_+^{N-l+k} \lambda_-^{l-k} + \lambda_+^{l-k} \lambda_-^{N-l+k}.
\label{eq:trace_corr_reorganizado}
\end{equation}
%
O denominador em \eqref{eq:traco_geral} é simplesmente
%
\begin{equation}
\operatorname{Tr}\left( T^N \right) = \lambda_+^N + \lambda_-^N.
\end{equation}
%
Substituindo \eqref{eq:trace_corr_reorganizado} e o denominador na expressão \eqref{eq:traco_geral}, obtemos finalmente
%
\begin{equation}
\boxed{
\langle \sigma_k \sigma_l \rangle
= \frac{ \lambda_+^{N-l+k} \lambda_-^{l-k} + \lambda_+^{l-k} \lambda_-^{N-l+k} }{ \lambda_+^N + \lambda_-^N }.
}
\label{eq:corr_final_passos}
\end{equation}
%
Fatorando $\lambda_+^N$ no numerador e denominador, obtém-se a forma que depende apenas da distância $|l-k|$ no limite termodinâmico, como discutido anteriormente.
%
\begin{equation}
\langle \sigma_k \sigma_l \rangle
= \frac{ \lambda_+^{N-l+k} \lambda_-^{l-k} + \lambda_+^{l-k} \lambda_-^{N-l+k} }{ \lambda_+^N + \lambda_-^N }.
\end{equation}
%
Fatorando $\lambda_+^N$ no numerador e denominador:
%
\begin{equation}
\langle \sigma_k \sigma_l \rangle
= \frac{ \left( \frac{\lambda_-}{\lambda_+} \right)^{l-k} + \left( \frac{\lambda_-}{\lambda_+} \right)^{N-l+k} }
{ 1 + \left( \frac{\lambda_-}{\lambda_+} \right)^N }.
\end{equation}
%
No limite termodinâmico ($N \to \infty$), o termo $\left( \frac{\lambda_-}{\lambda_+} \right)^{N-l+k}$ tende a zero, e o denominador tende a 1. Portanto,
%
\begin{equation}
\boxed{
\lim_{N \to \infty} \langle \sigma_k \sigma_l \rangle
= \left( \frac{\lambda_-}{\lambda_+} \right)^{l-k}
= \left( \frac{\lambda_-}{\lambda_+} \right)^{|l-k|}
}.
\label{eq:correlacao_geral}
\end{equation}
%
onde substituímos $l-k$ por $|l-k|$ para tornar a expressão válida para qualquer par de sítios.
%
Para $H = 0$ (isto é, $h = 0$), usando (\ref{eq:autovalores_h0}):
%
\begin{equation}
\frac{\lambda_-}{\lambda_+}
= \frac{2\sinh K}{2\cosh K}
= \tanh K.
\end{equation}
%
Assim, a função de correlação de pares em campo nulo é
%
\begin{equation}
\boxed{
\langle \sigma_k \sigma_l \rangle_{H=0}
= \left( \tanh K \right)^{|l-k|}
= \left[ \tanh\left( \frac{J}{k_B T} \right) \right]^{|l-k|}
}.
\label{eq:correlacao_h0}
\end{equation}
%
Alternativamente, podemos escrever
%
\begin{equation}
\langle \sigma_k \sigma_l \rangle_{H=0}
= \exp\left( - \frac{|l-k|}{\xi} \right),
\end{equation}
%
onde definimos o comprimento de correlação $\xi$ como
%
\begin{equation}
\boxed{
\frac{1}{\xi} = - \ln( \tanh K )
= - \ln\left[ \tanh\left( \frac{J}{k_B T} \right) \right]
}.
\label{eq:comprimento_correlacao}
\end{equation}
%
Ou ainda, na forma equivalente,
%
\begin{equation}
\boxed{
\xi = \frac{1}{ |\ln(\tanh K)| }.
}
\label{eq:xi_forma2}
\end{equation}
%
\subsection*{(b) Relação de flutuação-dissipação no equilíbrio}
%
A função de partição é
%
\begin{equation}
Z_N(T,H) = \sum_{\{\sigma_i\}} e^{-\beta \mathcal H},
\end{equation}
%
onde $\mathcal H = -J\sum_{i=1}^N \sigma_i\sigma_{i+1} - H\sum_{i=1}^N\sigma_i$. A magnetização total para uma configuração específica é $M = \sum_i \sigma_i$. Observamos que
%
\begin{equation}
\frac{\partial \mathcal H}{\partial H} = - \sum_{i=1}^N \sigma_i = -M.
\end{equation}
%
A média térmica da magnetização total é, por definição,
%
\begin{equation}
\langle M \rangle = \frac{1}{Z_N} \sum_{\{\sigma_i\}} M e^{-\beta \mathcal H}.
\end{equation}
%
Vamos calcular a derivada da função de partição em relação a $H$:
%
\begin{align}
\frac{\partial Z_N}{\partial H}
&= \sum_{\{\sigma_i\}} \frac{\partial}{\partial H} e^{-\beta \mathcal H} \\
&= \sum_{\{\sigma_i\}} \left( -\beta \frac{\partial \mathcal H}{\partial H} \right) e^{-\beta \mathcal H} \\
&= \sum_{\{\sigma_i\}} \beta M e^{-\beta \mathcal H} \\
&= \beta \, \langle M \rangle Z_N.
\label{eq:dZ_dH}
\end{align}
%
Portanto, isolando $\langle M \rangle$:
%
\begin{equation}
\langle M \rangle = \frac{1}{\beta} \frac{1}{Z_N} \frac{\partial Z_N}{\partial H}
= \frac{1}{\beta} \frac{\partial \ln Z_N}{\partial H}.
\label{eq:M_beta_lnZ}
\end{equation}
%
Esta é a relação usada no enunciado. Para obter a variância, precisamos da derivada segunda de $Z_N$ em relação a $H$. Derivando $\frac{\partial Z_N}{\partial H} = \beta \langle M \rangle Z_N$ novamente em relação a $H$, podemos proceder diretamente pela definição:
%
\begin{align}
\frac{\partial^2 Z_N}{\partial H^2}
&= \sum_{\{\sigma_i\}} \frac{\partial}{\partial H} \left( \beta M e^{-\beta \mathcal H} \right) \\
&= \sum_{\{\sigma_i\}} \beta M \left( -\beta \frac{\partial \mathcal H}{\partial H} \right) e^{-\beta \mathcal H} \\
&= \sum_{\{\sigma_i\}} \beta^2 M^2 e^{-\beta \mathcal H} \\
&= \beta^2 \, \langle M^2 \rangle Z_N.
\label{eq:d2Z_dH2}
\end{align}
%
Onde ao derivar $\beta M$ em relação a $H$, o fator $M$ é constante para uma configuração fixa (ele não depende de $H$ explicitamente, apenas a probabilidade de ocorrer muda), portanto a derivada de $\beta M$ é zero, restando apenas a derivada da exponencial.
%
Usando a equação \eqref{eq:M_beta_lnZ}, derivamos $\langle M \rangle$ mais uma vez em relação a $H$:
%
\begin{align}
\frac{\partial \langle M \rangle}{\partial H}
&= \frac{1}{\beta} \frac{\partial^2 \ln Z_N}{\partial H^2}.
\end{align}
%
Agora, calculamos a derivada segunda do logaritmo:
%
\begin{align}
\frac{\partial^2 \ln Z_N}{\partial H^2}
&= \frac{\partial}{\partial H} \left( \frac{1}{Z_N} \frac{\partial Z_N}{\partial H} \right) \\
&= \frac{1}{Z_N} \frac{\partial^2 Z_N}{\partial H^2}
- \frac{1}{Z_N^2} \left( \frac{\partial Z_N}{\partial H} \right)^2.
\label{eq:segunda_lnZ}
\end{align}
%
Substituímos as expressões \eqref{eq:dZ_dH} e \eqref{eq:d2Z_dH2} em \eqref{eq:segunda_lnZ}:
%
\begin{align}
\frac{\partial^2 \ln Z_N}{\partial H^2}
&= \frac{1}{Z_N} \left( \beta^2 \langle M^2 \rangle Z_N \right)
- \frac{1}{Z_N^2} \left( \beta \langle M \rangle Z_N \right)^2 \\
&= \beta^2 \langle M^2 \rangle - \beta^2 \langle M \rangle^2 \\
&= \beta^2 \left[ \langle M^2 \rangle - \langle M \rangle^2 \right].
\end{align}
%
Portanto,
%
\begin{align}
\frac{\partial \langle M \rangle}{\partial H}
&= \frac{1}{\beta} \left[ \beta^2 \left( \langle M^2 \rangle - \langle M \rangle^2 \right) \right] \\
&= \beta \left[ \langle M^2 \rangle - \langle M \rangle^2 \right].
\label{eq:dM_dH_final}
\end{align}
%
A magnetização por sítio é $m \equiv \frac{\langle M \rangle}{N}$. A susceptibilidade magnética é definida por
%
\begin{equation}
\chi(T,H) \equiv \left. \frac{\partial m}{\partial H} \right|_T
= \frac{1}{N} \frac{\partial \langle M \rangle}{\partial H}.
\end{equation}
%
Substituindo \eqref{eq:dM_dH_final}:
%
\begin{equation}
\chi(T,H) = \frac{\beta}{N} \left[ \langle M^2 \rangle - \langle M \rangle^2 \right].
\label{eq:chi_variance}
\end{equation}
%
Agora, expandimos os momentos da magnetização total em termos das variáveis de spin individuais. Temos
%
\begin{equation}
M = \sum_{i=1}^N \sigma_i,
\end{equation}
%
de modo que
%
\begin{equation}
M^2 = \sum_{i=1}^N \sum_{j=1}^N \sigma_i \sigma_j.
\end{equation}
%
Tomando a média térmica:
%
\begin{equation}
\langle M^2 \rangle = \sum_{i,j=1}^N \langle \sigma_i \sigma_j \rangle.
\end{equation}
%
Analogamente, para o quadrado da média:
%
\begin{equation}
\langle M \rangle^2 = \left( \sum_{i=1}^N \langle \sigma_i \rangle \right)^2
= \sum_{i,j=1}^N \langle \sigma_i \rangle \langle \sigma_j \rangle.
\end{equation}
%
Substituindo essas duas expansões na equação \eqref{eq:chi_variance}:
%
\begin{align}
\chi(T,H)
&= \frac{\beta}{N} \left[
\sum_{i,j=1}^N \langle \sigma_i \sigma_j \rangle
-
\sum_{i,j=1}^N \langle \sigma_i \rangle \langle \sigma_j \rangle
\right] \\
&= \frac{\beta}{N} \sum_{i,j=1}^N \left[
\langle \sigma_i \sigma_j \rangle - \langle \sigma_i \rangle \langle \sigma_j \rangle
\right].
\end{align}
%
O termo entre colchetes é a \textbf{função de correlação conectada} (ou cumulante) de dois pontos. Portanto, obtemos a relação de flutuação-dissipação na forma solicitada:
%
\begin{equation}
\boxed{
\chi(T,H)
= \frac{\beta}{N} \sum_{i,j=1}^N \left[
\langle \sigma_i \sigma_j \rangle - \langle \sigma_i \rangle \langle \sigma_j \rangle
\right].
}
\label{eq:fluc_diss_expandida}
\end{equation}
%
A equação \eqref{eq:fluc_diss_expandida} mostra que a susceptibilidade magnética (resposta a um campo externo) é proporcional à soma de todas as correlações de pares. 
%
\subsection*{(c) Susceptibilidade a campo nulo}
%
Em $H=0$ vale $\avg{\sigma_i}=0$ para todo $i$ (por simetria $\sigma\to-\sigma$ do hamiltoniano em campo nulo), de modo que
%
\begin{equation}
    \chi_0\equiv\chi(T,H\to0)=\frac\beta N\sum_{i,j}\avg{\sigma_i\sigma_j}.
\end{equation}
%
Em um sistema com simetria translacional (invariante por translação), as médias de um único sítio são todas iguais: $\langle \sigma_i \rangle = m$ para todo $i$. Além disso, a correlação entre dois sítios depende apenas da distância $r = |i-j|$. Nesse caso, podemos escrever
%
\begin{equation}
\sum_{i,j=1}^N \left[ \langle \sigma_i \sigma_j \rangle - \langle \sigma_i \rangle \langle \sigma_j \rangle \right]
= N \sum_{r=-\infty}^{\infty} \left[ \langle \sigma_0 \sigma_r \rangle - m^2 \right],
\end{equation}
%
de modo que a susceptibilidade por sítio se reduz a
%
\begin{equation}
\chi = \beta \sum_{r} \left[ \langle \sigma_0 \sigma_r \rangle - m^2 \right],
\end{equation}
%
usando invariância translacional, $\avg{\sigma_i\sigma_j}$ depende apenas de $r=j-i$, e a soma dupla se reduz a $N$ vezes a soma sobre $r$:
%
\begin{equation}
    \chi_0=\beta\sum_{r=-\infty}^{\infty}\avg{\sigma_0\sigma_r}
    =\beta\left[1+2\sum_{r=1}^\infty\left(\tanh\beta J\right)^{r}\right],
\end{equation}
%
onde usamos o resultado \eqref{eq:correlacao_h0} do item (a) e o limite termodinâmico $N\to\infty$. A série geométrica com razão $x=\tanh\beta J<1$ fornece $\sum_{r=1}^\infty x^r=x/(1-x)$, logo
%
\begin{equation}
    \chi_0=\beta\,\frac{1+\tanh\beta J}{1-\tanh\beta J}.
\end{equation}
%
Usando a identidade $\dfrac{1+\tanh y}{1-\tanh y}=\e{2y}$, obtemos finalmente
%
\begin{equation}
    \boxed{
    \chi_0(T)=\beta\,\e{2\beta J}=\frac{1}{k_BT}\exp\left(\frac{2J}{k_BT}\right).
    }
    \label{eq:chi0}
\end{equation}
%
Esse é o célebre resultado da cadeia de Ising 1D: não há transição de fase em $T>0$, mas $\chi_0$ diverge \emph{exponencialmente} (não como lei de potência) quando $T\to0$ -- o ponto crítico é empurrado exatamente para $T=0$. Para $T\to\infty$, $\chi_0\to\beta=1/(k_BT)$, recuperando a lei de Curie de spins não interagentes. A Fig.~\ref{fig:chi0} mostra o comportamento de $\chi_0(T)$.
%
\begin{figure}[H]
    \centering
    \includegraphics[width=0.62\linewidth]{Figures/statistical_mechanics/chi0_ising1d.png}
    \caption{Susceptibilidade a campo nulo da cadeia de Ising 1D, Eq.~\eqref{eq:chi0} (unidades $J=k_B=1$). $\chi_0$ diverge exponencialmente quando $T\to0$ e decai para a lei de Curie $\chi_0\sim1/T$ quando $T\to\infty$.}
    \label{fig:chi0}
\end{figure}
% ============================================================
\section*{Problema (2) -- Cadeia de spin 1 (modelo de Blume-Capel unidimensional)}
%
O hamiltoniano é
%
\begin{equation}
    \mathcal H=-J\sum_{i=1}^N S_iS_{i+1}+D\sum_{i=1}^N S_i^2,
    \qquad J>0,\ D>0,\ S_i=-1,0,+1,
    \label{eq:H_spin1}
\end{equation}
%
com condições periódicas de contorno. Primeiro, vamos definir as variáveis
%
\begin{equation}
    t\equiv\frac{k_BT}{J},
    \qquad
    d\equiv\frac{D}{J},
    \qquad
    K\equiv\beta J=\frac1t.
\end{equation}
%
Onde $d$ é o chamado parâmetro de campo cristalino.
%
\subsection*{(i) Autovalores da matriz de transferência}
%
Multiplicando o Hamiltoniano \eqref{eq:H_spin1} por \(\beta\), obtemos
%
\begin{equation}
    \beta \mathcal H = -K \sum_{i=1}^N S_i S_{i+1} + Kd \sum_{i=1}^N S_i^2.
\end{equation}
%
O termo de campo cristalino \(Kd S_i^2\) é local, dependendo de um único sítio \(i\). Para que a função de partição possa ser escrita como um produto de fatores associados a cada elo \((i, i+1)\), precisamos reescrevê-lo de forma simétrica entre os dois elos que incidem sobre o sítio \(i\): o elo \((i-1, i)\) e o elo \((i, i+1)\). 
%
Observamos a identidade
%
\begin{equation}
    \sum_{i=1}^N S_i^2 = \frac{1}{2} \sum_{i=1}^N \left( S_i^2 + S_{i+1}^2 \right),
    \label{eq:identidade}
\end{equation}
%
que é válida com condições periódicas de contorno, pois cada sítio aparece duas vezes na soma do lado direito (uma vez como \(S_i^2\) e outra como \(S_{i+1}^2\)), e o fator \(1/2\) compensa essa dupla contagem. Portanto,
%
\begin{equation}
    D \sum_{i=1}^N S_i^2 = \frac{D}{2} \sum_{i=1}^N \left( S_i^2 + S_{i+1}^2 \right).
\end{equation}
%
Multiplicando por \(\beta\) e usando \(Kd = \beta D\):
%
\begin{equation}
    \beta D \sum_{i=1}^N S_i^2 = \frac{Kd}{2} \sum_{i=1}^N \left( S_i^2 + S_{i+1}^2 \right).
\end{equation}
%
A função de partição canônica é
%
\begin{equation}
    Z_N = \sum_{\{S_i\}} e^{-\beta \mathcal H}.
\end{equation}
%
Substituindo a forma reescrita de \(\beta \mathcal H\):
%
\begin{equation}
    Z_N = \sum_{\{S_i\}} \exp\left\{ \sum_{i=1}^N \left[ K S_i S_{i+1} - \frac{Kd}{2} \left( S_i^2 + S_{i+1}^2 \right) \right] \right\}.
    \label{eq:Z_fatorada}
\end{equation}
%
Agora, o termo dentro da soma sobre \(i\) depende \textbf{apenas} dos dois spins consecutivos \(S_i\) e \(S_{i+1}\). Isso permite escrever a exponencial como um produto de fatores associados a cada elo:
%
\begin{equation}
    Z_N = \sum_{S_1, \dots, S_N} \prod_{i=1}^N T_{S_i, S_{i+1}},
\end{equation}
%
onde definimos o elemento de matriz da matriz de transferência \(T\) como
%
\begin{equation}
    T_{S, S'} \equiv \exp\left[ K S S' - \frac{Kd}{2} \left( S^2 + S'^2 \right) \right].
    \label{eq:T_elemento}
\end{equation}
%
Os spins podem assumir três valores: \(-1\), \(0\), e \(+1\). Escolhemos a base ordenada
%
\begin{equation}
    \left\{ -1,\; 0,\; +1 \right\}.
\end{equation}
%
Vamos calcular cada elemento \(T_{S,S'}\) usando a definição \eqref{eq:T_elemento} e as seguintes regras:
%
\begin{itemize}
    \item \(S^2 = 1\) se \(S = \pm 1\), e \(S^2 = 0\) se \(S = 0\);
    \item \(S S' = 1\) se \(S = S' = \pm 1\) (ambos \(+1\) ou ambos \(-1\));
    \item \(S S' = -1\) se \(S = -S'\) com \(S,S' = \pm 1\);
    \item \(S S' = 0\) se algum dos spins for \(0\).
\end{itemize}
%
Então os elementos de matriz (9 no total) serão,
%
\begin{equation}
    T_{1,1} = \exp\left[ K (1)(1) - \frac{Kd}{2} (1 + 1) \right]
    = \exp\left[ K - Kd \right]
    = e^{K(1-d)} \equiv a.
\end{equation}
%
\begin{equation}
    T_{-1,-1} = \exp\left[ K (-1)(-1) - \frac{Kd}{2} (1 + 1) \right]
    = \exp\left[ K - Kd \right]
    = e^{K(1-d)} = a.
\end{equation}
%
\begin{equation}
    T_{1,-1} = \exp\left[ K (1)(-1) - \frac{Kd}{2} (1 + 1) \right]
    = \exp\left[ -K - Kd \right]
    = e^{-K(1+d)} \equiv c.
\end{equation}
%
\begin{equation}
    T_{-1,1} = \exp\left[ K (-1)(1) - \frac{Kd}{2} (1 + 1) \right]
    = \exp\left[ -K - Kd \right]
    = e^{-K(1+d)} = c.
\end{equation}
%
\begin{equation}
    T_{1,0} = \exp\left[ K (1)(0) - \frac{Kd}{2} (1 + 0) \right]
    = \exp\left[ -\frac{Kd}{2} \right]
    \equiv b.
\end{equation}
%
\begin{equation}
    T_{0,1} = \exp\left[ K (0)(1) - \frac{Kd}{2} (0 + 1) \right]
    = \exp\left[ -\frac{Kd}{2} \right]
    = b.
\end{equation}
%
\begin{equation}
    T_{-1,0} = \exp\left[ K (-1)(0) - \frac{Kd}{2} (1 + 0) \right]
    = \exp\left[ -\frac{Kd}{2} \right]
    = b.
\end{equation}
%
\begin{equation}
    T_{0,-1} = \exp\left[ K (0)(-1) - \frac{Kd}{2} (0 + 1) \right]
    = \exp\left[ -\frac{Kd}{2} \right]
    = b.
\end{equation}
%
\begin{equation}
    T_{0,0} = \exp\left[ K (0)(0) - \frac{Kd}{2} (0 + 0) \right]
    = e^0 = 1.
\end{equation}
%
Organizando os elementos na base \(\{-1,0,+1\}\), obtemos a matriz
%
\begin{equation}
    T =
    \begin{pmatrix}
        T_{-1,-1} & T_{-1,0} & T_{-1,+1} \\
        T_{0,-1}  & T_{0,0}  & T_{0,+1}  \\
        T_{+1,-1} & T_{+1,0} & T_{+1,+1}
    \end{pmatrix}
    =
    \begin{pmatrix}
        a & b & c \\
        b & 1 & b \\
        c & b & a
    \end{pmatrix},
    \label{eq:T_final}
\end{equation}
%
onde definimos
%
\begin{equation}
    a \equiv e^{K(1-d)},
    \qquad
    b \equiv e^{-Kd/2},
    \qquad
    c \equiv e^{-K(1+d)}.
    \label{eq:def_abc}
\end{equation}
%
Com as condições periódicas \(S_{N+1} = S_1\), a soma sobre todas as configurações na equação \eqref{eq:Z_fatorada} equivale a tomar o traço da \(N\)-ésima potência da matriz de transferência:
%
\begin{equation}
    Z_N = \sum_{S_1, \dots, S_N} T_{S_1,S_2} T_{S_2,S_3} \cdots T_{S_N,S_1}
    = \operatorname{Tr}\left( T^N \right).
    \label{eq:Z_traco}
\end{equation}
%
Essa matriz é uma matriz supersimétrica podendo ser facilmente diagonalizada. A matriz~\eqref{eq:T_final} é invariante sob a troca simultânea das linhas e colunas $1\leftrightarrow3$ (simetria $S\to-S$). Isso permite bloco-diagonalizá-la nas combinações simétrica/antissimétrica e todas as quantidades termodinâmicas podem ser obtidas a partir dos autovalores de \(T\). Vamos bloco-diagonalizar \(T\) usando os seguintes vetores ortonormais:
%
\begin{equation}
    \ket{+} = \frac{1}{\sqrt{2}}
    \begin{pmatrix} 1 \\ 0 \\ 1 \end{pmatrix},
    \qquad
    \ket{0} = 
    \begin{pmatrix} 0 \\ 1 \\ 0 \end{pmatrix},
    \qquad
    \ket{-} = \frac{1}{\sqrt{2}}
    \begin{pmatrix} 1 \\ 0 \\ -1 \end{pmatrix}.
\label{eq:autovetores}    
\end{equation}
%
O setor antissimétrico é gerado por \(\ket{-}\), enquanto o setor simétrico é gerado por \(\{\ket{+}, \ket{0}\}\).
%
Vamos verificar que \(\ket{-}\) é um autovetor de \(T\):
%
\begin{equation}
    T \ket{-}
    =
    \frac{1}{\sqrt{2}}
    \begin{pmatrix}
        a & b & c \\
        b & 1 & b \\
        c & b & a
    \end{pmatrix}
    \begin{pmatrix} 1 \\ 0 \\ -1 \end{pmatrix}
    =
    \frac{1}{\sqrt{2}}
    \begin{pmatrix}
        a - c \\
        b - b \\
        c - a
    \end{pmatrix}
    =
    (a-c)
    \frac{1}{\sqrt{2}}
    \begin{pmatrix} 1 \\ 0 \\ -1 \end{pmatrix}.
\end{equation}
%
Portanto,
%
\begin{equation}
    \boxed{T \ket{-} = (a - c) \ket{-}},
    \qquad \Longrightarrow \qquad
    \boxed{\lambda_3 = a - c}.
\end{equation}
%
No subespaço \(\{\ket{+}, \ket{0}\}\), a matriz de transferência se restringe a um bloco \(2 \times 2\). Calculamos as imagens:
%
\begin{equation}
    T \ket{+}
    =
    \frac{1}{\sqrt{2}}
    \begin{pmatrix}
        a & b & c \\
        b & 1 & b \\
        c & b & a
    \end{pmatrix}
    \begin{pmatrix} 1 \\ 0 \\ 1 \end{pmatrix}
    =
    \frac{1}{\sqrt{2}}
    \begin{pmatrix}
        a + c \\
        2b \\
        a + c
    \end{pmatrix}
    =
    (a+c) \ket{+} + \sqrt{2} b \ket{0}.
\end{equation}
%
\begin{equation}
    T \ket{0}
    =
    \begin{pmatrix}
        a & b & c \\
        b & 1 & b \\
        c & b & a
    \end{pmatrix}
    \begin{pmatrix} 0 \\ 1 \\ 0 \end{pmatrix}
    =
    \begin{pmatrix}
        b \\
        1 \\
        b
    \end{pmatrix}
    =
    \sqrt{2} b \ket{+} + 1 \cdot \ket{0}.
\end{equation}
%
Assim, a matriz do bloco simétrico na base \(\{\ket{+}, \ket{0}\}\) é
%
\begin{equation}
    M_{\text{sim}} =
    \begin{pmatrix}
        a + c & \sqrt{2}\, b \\
        \sqrt{2}\, b & 1
    \end{pmatrix}.
    \label{eq:M_sim}
\end{equation}
%
Os autovalores \(\lambda\) de \(M_{\text{sim}}\) são as raízes da equação secular
%
\begin{equation}
    \det\left( M_{\text{sim}} - \lambda I \right) = 0.
\end{equation}
%
Substituindo a matriz \eqref{eq:M_sim}:
%
\begin{equation}
    \det
    \begin{pmatrix}
        a + c - \lambda & \sqrt{2}\, b \\
        \sqrt{2}\, b & 1 - \lambda
    \end{pmatrix}
    = 0.
\end{equation}
%
Calculando o determinante:
%
\begin{equation}
    (a + c - \lambda)(1 - \lambda) - 2b^2 = 0.
\end{equation}
%
Expandindo:
%
\begin{equation}
    (a + c - \lambda)(1 - \lambda) - 2b^2 = 0,
\end{equation}
%
\begin{equation}
    (a + c) - (a + c)\lambda - \lambda + \lambda^2 - 2b^2 = 0,
\end{equation}
%
\begin{equation}
    \lambda^2 - (a + c + 1)\lambda + (a + c) - 2b^2 = 0.
\end{equation}
%
Portanto, a equação secular é
%
\begin{equation}
    \boxed{
    \lambda^2 - (a + c + 1)\lambda + \left[ (a + c) - 2b^2 \right] = 0.
    }
    \label{eq:secular}
\end{equation}
%
Resolvendo a equação quadrática \eqref{eq:secular}:
%
\begin{equation}
    \lambda_{\pm} = \frac{(a + c + 1) \pm \sqrt{(a + c + 1)^2 - 4\left[ (a + c) - 2b^2 \right]}}{2}.
\end{equation}
%
Simplificando o discriminante:
%
\begin{align}
    \Delta &= (a + c + 1)^2 - 4(a + c) + 8b^2 \\
    &= (a + c)^2 + 2(a + c) + 1 - 4(a + c) + 8b^2 \\
    &= (a + c)^2 - 2(a + c) + 1 + 8b^2 \\
    &= (a + c - 1)^2 + 8b^2.
\end{align}
%
Portanto,
%
\begin{equation}
    \boxed{
    \lambda_{\pm} = \frac{(a + c + 1) \pm \sqrt{(a + c - 1)^2 + 8b^2}}{2}.
    }
    \label{eq:autovalores_simetricos}
\end{equation}
%
Como \(a, b, c > 0\) para toda temperatura finita (\(K < \infty\)), o discriminante \(\Delta = (a+c-1)^2 + 8b^2 > 0\). Além disso, \(\sqrt{\Delta} > |a+c-1|\), de modo que
%
\begin{equation}
    \lambda_+ = \frac{(a + c + 1) + \sqrt{\Delta}}{2}
    \quad > \quad
    \lambda_- = \frac{(a + c + 1) - \sqrt{\Delta}}{2}.
\end{equation}
%
Reunindo os resultados:
%
\begin{itemize}
    \item Do setor simétrico:
    \begin{equation}
        \lambda_1 = \lambda_+ = \frac{(a + c + 1) + \sqrt{(a + c - 1)^2 + 8b^2}}{2},
    \end{equation}
    \begin{equation}
        \lambda_2 = \lambda_- = \frac{(a + c + 1) - \sqrt{(a + c - 1)^2 + 8b^2}}{2}.
    \end{equation}

    \item Do setor antissimétrico:
    \begin{equation}
        \lambda_3 = a - c.
    \end{equation}
\end{itemize}
%
Portanto, os três autovalores são
%
\begin{equation}
    \boxed{
    \begin{aligned}
        \lambda_1 &= \frac{(a + c + 1) + \sqrt{(a + c - 1)^2 + 8b^2}}{2}, \\[6pt]
        \lambda_2 &= \frac{(a + c + 1) - \sqrt{(a + c - 1)^2 + 8b^2}}{2}, \\[6pt]
        \lambda_3 &= a - c.
    \end{aligned}
    }
    \label{eq:autovalores_completos}
\end{equation}
%
Pelo teorema de Perron–Frobenius~\cite{MacCluer2000}, como todos os elementos de \(T\) são estritamente positivos, o maior autovalor é único, positivo, e possui um autovetor com componentes de mesmo sinal. O autovetor associado a \(\lambda_3 = a - c\) é \(\ket{-} = (1,0,-1)/\sqrt{2}\), cujas componentes não têm todas o mesmo sinal. Portanto, \(\lambda_3\) \textbf{não} pode ser o autovalor dominante. Resta comparar \(\lambda_1\) e \(\lambda_2\). Como \(\lambda_1 > \lambda_2\) (pois a raiz quadrada é positiva), concluímos que
%
\begin{equation}
    \boxed{
    \lambda_1 > \lambda_2,\ \lambda_3
    \qquad \text{para toda temperatura finita } (K < \infty).
    }
\end{equation}
%
Assim,
%
\begin{equation}
    \lambda_{\max} = \lambda_1 = \frac{(a + c + 1) + \sqrt{(a + c - 1)^2 + 8b^2}}{2}.
\end{equation}
%
\subsection*{(ii) Energia interna: valor esperado do hamiltoniano}
%
Pelo exercício anterior, sabemos que
%
\begin{equation}
    \lambda_1 > \lambda_2, \qquad \lambda_1 > |\lambda_3|,
\end{equation}
%
de modo que \(\lambda_1 = \lambda_{\max}\) é o maior autovalor em módulo (e positivo). Podemos escrever
%
\begin{equation}
    T = \lambda_1 \ket{1}\bra{1} + \lambda_2 \ket{2}\bra{2} + \lambda_3 \ket{3}\bra{3},
\end{equation}
%
onde \(\ket{i}\) são os autovetores ortonormais correspondentes definidos em~\eqref{eq:autovetores}. Elevando à potência \(N\):
%
\begin{equation}
    T^N = \lambda_1^N \ket{1}\bra{1} + \lambda_2^N \ket{2}\bra{2} + \lambda_3^N \ket{3}\bra{3}.
\end{equation}
%
Tomando o traço:
%
\begin{equation}
    Z_N = \operatorname{Tr}(T^N) = \lambda_1^N + \lambda_2^N + \lambda_3^N.
    \label{eq:Z_autovalores}
\end{equation}
%
Fatoramos o termo dominante \(\lambda_1^N\) na equação \eqref{eq:Z_autovalores}:
%
\begin{equation}
    Z_N = \lambda_1^N \left[ 1 + \left( \frac{\lambda_2}{\lambda_1} \right)^N + \left( \frac{\lambda_3}{\lambda_1} \right)^N \right].
    \label{eq:Z_fatorado}
\end{equation}
%
Como \(\lambda_1\) é estritamente maior que \(\lambda_2\) e \(|\lambda_3|\), temos
%
\begin{equation}
    \left| \frac{\lambda_2}{\lambda_1} \right| < 1,
    \qquad
    \left| \frac{\lambda_3}{\lambda_1} \right| < 1.
\end{equation}
%
Portanto, no limite \(N \to \infty\),
%
\begin{equation}
    \lim_{N \to \infty} \left( \frac{\lambda_2}{\lambda_1} \right)^N = 0,
    \qquad
    \lim_{N \to \infty} \left( \frac{\lambda_3}{\lambda_1} \right)^N = 0.
\end{equation}
%
Assim,
%
\begin{equation}
    \lim_{N \to \infty} \frac{Z_N}{\lambda_1^N}
    = \lim_{N \to \infty} \left[ 1 + \left( \frac{\lambda_2}{\lambda_1} \right)^N + \left( \frac{\lambda_3}{\lambda_1} \right)^N \right]
    = 1.
\end{equation}
%
Ou seja,
%
\begin{equation}
    Z_N \sim \lambda_1^N
    \qquad (N \to \infty),
    \label{eq:Z_approx}
\end{equation}
%
onde o símbolo \(\sim\) significa que a razão \(Z_N / \lambda_1^N\) tende a 1.
%
Tomando o logaritmo natural de \eqref{eq:Z_fatorado}:
%
\begin{equation}
    \ln Z_N = N \ln \lambda_1 + \ln\left[ 1 + \left( \frac{\lambda_2}{\lambda_1} \right)^N + \left( \frac{\lambda_3}{\lambda_1} \right)^N \right].
    \label{eq:lnZ_exato}
\end{equation}
%
Assim, no limite termodinâmico,
%
\begin{equation}
    \ln Z_N = N \ln \lambda_1 + o(N),
\end{equation}
%
onde \(o(N)\) denota um termo que cresce mais lentamente que \(N\) (neste caso, tende a uma constante ou a zero). Portanto,
%
\begin{equation}
    \ln Z_N \approx N \ln \lambda_1
    \qquad (N \to \infty).
    \label{eq:lnZ_approx}
\end{equation}
%
Agora, o valor esperado do Hamiltoniano (energia interna) é definido como
%
\begin{equation}
    U \equiv \langle \mathcal H \rangle
    = \frac{1}{Z_N} \sum_{\{S_i\}} \mathcal H(\{S_i\}) e^{-\beta \mathcal H(\{S_i\})}.
    \label{eq:U_def}
\end{equation}
%
Para mostrar esta relação, calculamos a derivada de \(Z_N\) em relação a \(\beta\):
%
\begin{align}
    \frac{\partial Z_N}{\partial \beta}
    &= \frac{\partial}{\partial \beta} \sum_{\{S_i\}} e^{-\beta \mathcal H(\{S_i\})} \\
    &= \sum_{\{S_i\}} \left( -\mathcal H(\{S_i\}) \right) e^{-\beta \mathcal H(\{S_i\})} \\
    &= - \sum_{\{S_i\}} \mathcal H(\{S_i\}) e^{-\beta \mathcal H(\{S_i\})}.
    \label{eq:dZ_dbeta}
\end{align}
%
Calculamos a derivada de \(\ln Z_N\) em relação a \(\beta\):
%
\begin{align}
    \frac{\partial \ln Z_N}{\partial \beta}
    &= \frac{1}{Z_N} \frac{\partial Z_N}{\partial \beta} \\
    &= \frac{1}{Z_N} \left( - \sum_{\{S_i\}} \mathcal H(\{S_i\}) e^{-\beta \mathcal H(\{S_i\})} \right) \\
    &= - \frac{1}{Z_N} \sum_{\{S_i\}} \mathcal H(\{S_i\}) e^{-\beta \mathcal H(\{S_i\})}.
    \label{eq:dlnZ_dbeta}
\end{align}
%
Comparando a equação \eqref{eq:dlnZ_dbeta} com a definição de \(U\) em \eqref{eq:U_def}, vemos que
%
\begin{equation}
    \frac{\partial \ln Z_N}{\partial \beta} = -U.
\end{equation}
%
Portanto,
%
\begin{equation}
    U = - \frac{\partial \ln Z_N}{\partial \beta}.
    \label{eq:U_beta}
\end{equation}
%
Então, no limite termodinâmica a a energia interna é dada por, 
%
\begin{equation}
    U=-\pd{\ln Z_N}{\beta}=-J\pd{\ln Z_N}{K}=-NJ\,\pd{\ln\lambda_1}{K},
\end{equation}
%
de modo que a energia interna por sítio, em unidades de $J$, é
%
\begin{equation}
    \frac{U}{NJ}=-\pd{\ln\lambda_1}{K}
    =-\frac{1}{\lambda_1}\pd{\lambda_1}{K}.
    \label{eq:U_spin1}
\end{equation}
%
Podemos definir o maior autovalor de \(T\) como,
%
\begin{equation}
    \lambda_1 = \frac{S + 1 + R}{2},
\end{equation}
%
com as definições auxiliares
%
\begin{equation}
    S \equiv a + c,
    \qquad
    R \equiv \sqrt{(S - 1)^2 + 8b^2}.
\end{equation}
%
Começamos derivando \(\lambda_1\):
%
\begin{equation}
    \frac{\partial \lambda_1}{\partial K}
    = \frac{1}{2} \left( \frac{\partial S}{\partial K} + \frac{\partial R}{\partial K} \right).
    \label{eq:dlamb1}
\end{equation}
%
Da definição de \(R\), temos utilizando a regra da cadeia:
%
\begin{align}
    \frac{\partial R}{\partial K}
    &= \frac{1}{2R} \left[ 2(S - 1) \frac{\partial S}{\partial K} + 16 b \frac{\partial b}{\partial K} \right] \\
    &= \frac{1}{R} \left[ (S - 1) \frac{\partial S}{\partial K} + 8 b \frac{\partial b}{\partial K} \right].
    \label{eq:dR}
\end{align}
%
Sabemos que
%
\[
b = e^{-Kd/2}.
\]
%
Derivando em relação a \(K\):
%
\begin{equation}
    \frac{\partial b}{\partial K} = -\frac{d}{2} e^{-Kd/2} = -\frac{d}{2} b.
    \label{eq:db}
\end{equation}
%
Substituindo \eqref{eq:db} em \eqref{eq:dR}:
%
\begin{equation}
    8 b \frac{\partial b}{\partial K}
    = 8 b \left( -\frac{d}{2} b \right)
    = -4 d b^2.
\end{equation}
%
Portanto,
%
\begin{equation}
    \frac{\partial R}{\partial K}
    = \frac{1}{R} \left[ (S - 1) \frac{\partial S}{\partial K} - 4 d b^2 \right].
    \label{eq:dR_final}
\end{equation}
%
Como \(S = a + c\), temos
%
\begin{equation}
    \frac{\partial S}{\partial K} = \frac{\partial a}{\partial K} + \frac{\partial c}{\partial K}.
\end{equation}
%
Utilizando as definições~\eqref{eq:def_abc}, as derivadas de \(a\) e \(c\) são:
%
\[
\frac{\partial a}{\partial K} = (1-d) a,
\qquad
\frac{\partial c}{\partial K} = -(1+d) c.
\]
%
Portanto,
%
\begin{equation}
    \frac{\partial S}{\partial K} = (1-d) a - (1+d) c.
    \label{eq:dS}
\end{equation}
%
Substituindo \eqref{eq:dR_final} e \eqref{eq:dS} em \eqref{eq:dlamb1}:
%
\begin{align}
    \frac{\partial \lambda_1}{\partial K}
    &= \frac{1}{2} \left[
        \frac{\partial S}{\partial K}
        + \frac{1}{R} \left( (S - 1) \frac{\partial S}{\partial K} - 4 d b^2 \right)
    \right] \\
    &= \frac{1}{2} \left[
        \frac{\partial S}{\partial K}
        + \frac{(S - 1) \frac{\partial S}{\partial K} - 4 d b^2}{R}
    \right].
\label{eq:lambda_1_partial}    
\end{align}
%
Podemos escrever tudo explicitamente em função de \(K\) e \(d\). Temos:
%
\begin{align}
    a &= e^{K(1-d)}, \\
    c &= e^{-K(1+d)}, \\
    b &= e^{-Kd/2}, \\
    S &= e^{K(1-d)} + e^{-K(1+d)}, \\
    R &= \sqrt{ \left( e^{K(1-d)} + e^{-K(1+d)} - 1 \right)^2 + 8 e^{-Kd} }, \\
    \frac{\partial S}{\partial K} &= (1-d) e^{K(1-d)} - (1+d) e^{-K(1+d)}.
\end{align}
%
Substituindo esses valores em \eqref{eq:lambda_1_partial}, obtemos a forma totalmente explícita:
%
\begin{align}
    \frac{\partial \lambda_1}{\partial K}
    &= \frac{1}{2} \Bigg[
        (1-d) e^{K(1-d)} - (1+d) e^{-K(1+d)} \nonumber \\
        &\quad + \frac{
            \left( e^{K(1-d)} + e^{-K(1+d)} - 1 \right)
            \left[ (1-d) e^{K(1-d)} - (1+d) e^{-K(1+d)} \right]
            - 4d e^{-Kd}
        }{
            \sqrt{ \left( e^{K(1-d)} + e^{-K(1+d)} - 1 \right)^2 + 8 e^{-Kd} }
        }
    \Bigg].
    \label{eq:dlamb1_explicita}
\end{align}
%
Usando \(\lambda_1 = (S + 1 + R)/2\), temos
%
\[
\frac{1}{\lambda_1} = \frac{2}{S + 1 + R}.
\]
%
Portanto,
%
\begin{align}
    \frac{U}{NJ}
    &= - \frac{1}{\lambda_1} \frac{\partial \lambda_1}{\partial K} \\
    &= - \frac{2}{S + 1 + R}
    \cdot \frac{1}{2} \left[
        \frac{\partial S}{\partial K}
        + \frac{(S - 1) \frac{\partial S}{\partial K} - 4 d b^2}{R}
    \right] \\
    &= - \frac{1}{S + 1 + R}
    \left[
        \frac{\partial S}{\partial K}
        + \frac{(S - 1) \frac{\partial S}{\partial K} - 4 d b^2}{R}
    \right].
\end{align}
%
Ou seja,
%
\begin{equation}
    \boxed{
    \frac{U}{NJ} =
    - \frac{1}{S + 1 + R}
    \left[
        \frac{\partial S}{\partial K}
        + \frac{(S - 1) \frac{\partial S}{\partial K} - 4 d b^2}{R}
    \right]
    }
    \label{eq:U_explicita}
\end{equation}
%
com
%
\[
S = a + c, \qquad R = \sqrt{(S - 1)^2 + 8b^2}, \qquad \frac{\partial S}{\partial K} = (1-d)a - (1+d)c.
\]
%
Para verificar a consistência, calculamos o limite \(K \to 0\). Neste limite:
%
\[
a \to 1, \qquad b \to 1, \qquad c \to 1,
\]
\[
S \to 2, \qquad R \to \sqrt{(2-1)^2 + 8} = 3,
\]
\[
\frac{\partial S}{\partial K} \to (1-d) - (1+d) = -2d.
\]
%
Substituindo em \eqref{eq:U_explicita}:
%
\[
\frac{U}{NJ}
= - \frac{1}{2 + 1 + 3}
\left[
    -2d + \frac{(2-1)(-2d) - 4d(1)^2}{3}
\right].
\]
%
Simplificando o termo dentro dos colchetes:
%
\[
-2d + \frac{-2d - 4d}{3}
= -2d + \frac{-6d}{3}
= -2d - 2d = -4d.
\]
%
Portanto,
%
\[
\frac{U}{NJ} = - \frac{1}{6} (-4d) = \frac{2d}{3}.
\]
%
Este é o resultado esperado: em altas temperaturas, todos os três estados são igualmente prováveis, e
%
\[
\langle S_i^2 \rangle = \frac{0^2 + 1^2 + 1^2}{3} = \frac{2}{3},
\]
%
de modo que \(U/(NJ) = d \langle S_i^2 \rangle = 2d/3\).
%
Para \(K \to \infty\) e \(d < 1\), temos:
%
\[
a = e^{K(1-d)} \gg 1, \qquad b \to 0, \qquad c \to 0.
\]
%
Assim,
%
\[
S \approx a, \qquad R \approx a - 1, \qquad \lambda_1 \approx a,
\]
\[
\frac{\partial S}{\partial K} \approx (1-d)a.
\]
%
Substituindo em \eqref{eq:U_explicita}:
%
\[
\frac{U}{NJ}
\approx - \frac{1}{a + 1 + a}
\left[
    (1-d)a + \frac{(a-1)(1-d)a}{a}
\right]
= - \frac{1}{2a}
\left[ (1-d)a + (1-d)a \right]
= -(1-d) = d - 1.
\]
%
Portanto,
%
\[
\frac{U}{NJ} \to d - 1,
\]
%
que é a energia do estado fundamental ferromagnético, \(\varepsilon_F/J = -1 + d\).
%
Para \(K \to \infty\) e \(d > 1\), temos:
%
\[
a = e^{-K(d-1)} \to 0, \qquad c = e^{-K(d+1)} \to 0, \qquad b = e^{-Kd/2} \to 0,
\]
de modo que \(S \to 0\), \(R \to 1\), e \(\lambda_1 \to 1\).
%
Além disso,
%
\[
\frac{\partial S}{\partial K} \to 0.
\]
%
Substituindo em \eqref{eq:U_explicita}:
%
\[
\frac{U}{NJ}
\approx - \frac{1}{0 + 1 + 1}
\left[
    0 + \frac{(0-1) \cdot 0 - 4d b^2}{1}
\right]
= - \frac{1}{2} (-4d b^2) = 2d b^2.
\]
%
Mas \(b^2 = e^{-Kd}\), que tende a zero exponencialmente. Portanto,
%
\[
\frac{U}{NJ} \to 0,
\]
%
que é a energia do estado fundamental \(S_i = 0\).
%
Para \(d = 1\), temos \(a = e^{0} = 1\) e \(c = e^{-2K} \to 0\), \(b = e^{-K/2} \to 0\). Assim,
%
\[
S \to 1, \qquad R \to \sqrt{8} b = 2\sqrt{2} e^{-K/2}, \qquad \lambda_1 \approx 1 + \sqrt{2} e^{-K/2},
\]
\[
\frac{\partial S}{\partial K} = (0)a - 2c = -2e^{-2K} \to 0.
\]
%
Substituindo em \eqref{eq:U_explicita}:
%
\[
\frac{U}{NJ}
\approx - \frac{1}{1 + 1 + 2\sqrt{2} e^{-K/2}}
\left[
    0 + \frac{(0) \cdot 0 - 4(1)b^2}{2\sqrt{2} e^{-K/2}}
\right]
= - \frac{1}{2 + 2\sqrt{2} e^{-K/2}}
\left[ - \frac{4 e^{-K}}{2\sqrt{2} e^{-K/2}} \right].
\]
%
Simplificando:
%
\[
\frac{U}{NJ}
\approx \frac{1}{2} \cdot \frac{4 e^{-K}}{2\sqrt{2} e^{-K/2}}
= \frac{1}{\sqrt{2}} e^{-K/2}.
\]
%
Este resultado mostra que a energia interna tende a zero com uma correção \(\sim e^{-K/2}\), que é mais lenta do que a correção \(e^{-K}\) obtida para \(d > 1\). Isso reflete a degenerescência extra do estado fundamental no ponto \(d = 1\). A Fig.~\ref{fig:U_spin1} mostra $U/(NJ)$ em função de $t=k_BT/J$, obtido numericamente a partir da Eq.~\eqref{eq:U_spin1}, para vários valores de $d=D/J$: para $d<1$ a energia parte do valor fundamental ferromagnético $-1+d$ (ver item iii) e cresce monotonicamente até $\tfrac23 d$; para $d>1$ ela parte de $0$ (estado fundamental $S=0$) e também cresce monotonicamente até $\tfrac23 d$.
%
\begin{figure}[H]
    \centering
    \includegraphics[width=0.68\linewidth]{Figures/statistical_mechanics/energy_spin1.png}
    \caption{Energia interna por sítio $U/(NJ)$, Eq.~\eqref{eq:U_spin1}, em função da temperatura reduzida $t=k_BT/J$, para vários valores do parâmetro de campo cristalino $d=D/J$. Note a mudança de comportamento entre $d<1$ (fundamental ferromagnético, $U/NJ\to-1+d$ quando $t\to0$) e $d>1$ (fundamental $S=0$, $U/NJ\to0$ quando $t\to0$).}
    \label{fig:U_spin1}
\end{figure}
%
\subsection*{(iii) Estado fundamental e forma assintótica dos autovalores para $T\to0$}
%
No limite $T=0$ o sistema minimiza $\mathcal H$. Comparamos os dois candidatos naturais por sítio:
%
\begin{itemize}
    \item \textbf{Estado ferromagnético}, $S_i=+1$ (ou $-1$) para todo $i$: energia por sítio $\varepsilon_F/J=-1+d$.
    \item \textbf{Estado ``quadrupolar''}, $S_i=0$ para todo $i$: energia por sítio $\varepsilon_0/J=0$.
\end{itemize}
%
Comparando, $\varepsilon_F<\varepsilon_0 \iff -1+d<0\iff d<1$. Logo,
%
\begin{equation}
    \boxed{
    \text{Estado fundamental: }
    \begin{cases}
        \text{ferromagnético }(S_i=\pm1\text{ uniforme}), & d<1,\\[4pt]
        S_i=0\text{ uniforme (não-magnético)}, & d>1.
    \end{cases}
    }
\end{equation}
%
Em $d=1$ exatamente, $\varepsilon_F=\varepsilon_0=0$: esse é o célebre \emph{ponto de degenerescência} do modelo de Blume-Capel unidimensional a $T=0$, onde há uma degenerescência macroscópica adicional do estado fundamental (podem-se inserir domínios de $S=0$ dentro do estado ferromagnético sem custo de energia).
%
\medskip
%
\noindent\textbf{Região $d<1$ (fundamental ferromagnético).} Quando $K\to\infty$, $a=\e{K(1-d)}\to\infty$ domina exponencialmente sobre $b=\e{-Kd/2}\to0$ e $c=\e{-K(1+d)}\to0$. Expandindo $R=\sqrt{(a-1)^2+8b^2+\cdots}\approx a-1+4b^2/a$ para $a\gg1$, obtemos
%
\begin{align}
    \lambda_+&\approx a+\frac{2b^2}{a}=a+2\e{-K}=\e{K(1-d)}+2\e{-K}+\cdots,\\
    \lambda_0&\approx 1-\frac{2b^2}{a}=1-2\e{-K}+\cdots,\\
    \lambda_3&=a-c=\e{K(1-d)}-\e{-K(1+d)}\approx\e{K(1-d)}-\e{-K(1+d)}.
\end{align}
%
Portanto,
%
\begin{equation}
    \boxed{
    \lambda_1\sim\lambda_3\sim\e{K(1-d)}=\e{(1-d)/t}
    \qquad(t\to0^+,\ d<1),
    }
\end{equation}
%
consistente com $\varepsilon_F/J=-(1-d)$ (isto é, $\lambda_1\sim\e{-\beta N\varepsilon_F/N}=\e{-\beta\varepsilon_F}$). Além disso $\lambda_1-\lambda_3\approx2\e{-K}+\e{-K(1+d)}\to0$ exponencialmente: os dois autovalores tornam-se assintoticamente degenerados quando $T\to0$, refletindo a degenerescência dupla ($S=+1$ ou $S=-1$ uniforme) do estado fundamental ferromagnético -- a simetria só é restaurada dinamicamente no limite estrito $T\to0$, como esperado em uma dimensão (nenhuma quebra espontânea de simetria em $T>0$).
%
\medskip
\noindent\textbf{Região $d>1$ (fundamental $S=0$).} Agora, quando $K\to\infty$, $a=\e{-K(d-1)}\to0$, $c=\e{-K(d+1)}\to0$ ainda mais rápido, e $b=\e{-Kd/2}\to0$. Expandindo cuidadosamente até segunda ordem (os termos lineares em $a$ se cancelam),
%
\begin{equation}
    R=\sqrt{(a-1)^2+8b^2}\approx1-a+4b^2,
    \qquad
    \lambda_+=\frac{S+1+R}{2}\approx1+2b^2,
\end{equation}
%
de modo que
%
\begin{equation}
    \boxed{
    \lambda_+\approx1+2\e{-Kd}=1+2\e{-d/t}
    \qquad(t\to0^+,\ d>1),
    }
\end{equation}
%
consistente com $\varepsilon_0=0$ (fundamental não-degenerado, $\lambda_+\to1$), com uma correção ativada $\sim\e{-d/t}$ que reflete o ``gap'' de energia $\Delta\varepsilon\sim D$ para excitar um spin de $S=0$ para $S=\pm1$. Os outros dois autovalores, $\lambda_0\approx a-2b^2\sim\e{-K(d-1)}$ e $\lambda_3\approx a\sim\e{-K(d-1)}$, vão a zero, refletindo o gap para os estados excitados.
%
\medskip
\noindent\textbf{Ponto especial $d=1$.} Aqui $a=\e{K(1-d)}=1$ exatamente para todo $K$, enquanto $c=\e{-2K}\to0$ e $b=\e{-K/2}\to0$ mais devagar que $c$. Logo $S\approx1$, $R\approx\sqrt{8}\,b=2\sqrt2\,\e{-K/2}$, e
%
\begin{equation}
    \boxed{
    \lambda_+\approx1+\sqrt2\,\e{-K/2}=1+\sqrt2\,\e{-1/(2t)}
    \qquad(t\to0^+,\ d=1),
    }
\end{equation}
%
com correção $\sim\e{-1/(2t)}$, que decai \emph{mais devagar} que a correção $\e{-d/t}\big|_{d\to1}=\e{-1/t}$ obtida no limite $d\to1^+$ da região anterior. Isso é a assinatura, na matriz de transferência, da degenerescência extra do ponto $d=1$ mencionada acima: mais estados de baixa energia contribuem, dando uma aproximação mais lenta ao valor assintótico $\lambda_+\to1$~\cite{krinsky}.
% ============================================================
\section*{Problema (3) -- Modelo ANNNI unidimensional}
%
O hamiltoniano de spin com interações competitivas entre primeiros e segundos vizinhos é
%
\begin{equation}
    \mathcal H=-J\sum_{i=1}^N\sigma_i\sigma_{i+1}+pJ\sum_{i=1}^N\sigma_i\sigma_{i+2},
    \qquad J,p>0,
    \label{eq:H_annni}
\end{equation}
%
com interação ferromagnética entre primeiros vizinhos e antiferromagnética (competitiva) entre segundos vizinhos, como indicado no enunciado. Definimos $K\equiv\beta J$ e a temperatura reduzida $t=k_BT/J=1/K$.
%
\subsection*{(i) Estado fundamental em função do parâmetro de competição $p$}
%
O parâmetro \(p\) controla a intensidade relativa da competição.
%
No estado fundamental (\(T=0\)), o sistema minimiza a energia por sítio
%
\begin{equation}
    \varepsilon \equiv \frac{1}{N} \langle \mathcal H \rangle
    = -J \langle \sigma_i \sigma_{i+1} \rangle + pJ \langle \sigma_i \sigma_{i+2} \rangle.
\end{equation}
%
Para encontrar o estado fundamental, comparamos as energias de todas as configurações periódicas possíveis. Em uma cadeia unidimensional com interações até segunda vizinhança, os candidatos naturais são fases com períodos \(P = 1, 2, 3, 4, \dots\). Vamos analisar as principais.
%
\begin{itemize}
    \item \textbf{Ferromagnético,} todos os spins são iguais: \(\sigma_i = +1\) para todo \(i\) (ou todos \(-1\), que é equivalente por simetria de inversão global). Temos
%
\begin{equation}
    \langle \sigma_i \sigma_{i+1} \rangle = 1,
    \qquad
    \langle \sigma_i \sigma_{i+2} \rangle = 1.
\end{equation}
%
Portanto, a energia por sítio é
%
\begin{equation}
    \varepsilon_F = -J + pJ = J(-1 + p).
    \label{eq:eps_F}
\end{equation}   
%   
    \item \textbf{Fase antiferromagnética (AF),} a configuração alternada: \(\sigma_i = (-1)^i\) (ou \((-1)^{i+1}\)). Para esta fase:
%
\begin{equation}
    \sigma_i \sigma_{i+1} = (-1)^i (-1)^{i+1} = -1,
    \qquad
    \sigma_i \sigma_{i+2} = (-1)^i (-1)^{i+2} = +1.
\end{equation}
%
Logo,
%
\begin{equation}
    \langle \sigma_i \sigma_{i+1} \rangle = -1,
    \qquad
    \langle \sigma_i \sigma_{i+2} \rangle = +1,
\end{equation}
%
e a energia por sítio é
%
\begin{equation}
    \varepsilon_{AF} = -J(-1) + pJ(+1) = J(1 + p).
    \label{eq:eps_AF}
\end{equation}
%
    \item \textbf{Fase de período 4 (P4),} uma configuração com período 4, por exemplo, a sequência \( (+,+,-,-) \) repetida. Vamos calcular as médias para esta fase. A sequência é:
%
\[
\sigma: \quad +,\ +,\ -,\ -,\ +,\ +,\ -,\ -,\ \dots
\]
%
Para esta fase, calculamos: primeiros vizinhos: pares \((+,+)\), \((+,-)\), \((-,-)\), \((-,+)\). Média:
  \[
  \langle \sigma_i \sigma_{i+1} \rangle = \frac{1}{4} \left[ (+1)(+1) + (+1)(-1) + (-1)(-1) + (-1)(+1) \right]
  = \frac{1}{4} (1 - 1 + 1 - 1) = 0.
  \]
%
Segundos vizinhos: pares \((+, -)\), \((+ , -)\), \((- , +)\), \((- , +)\). Média:
  \[
  \langle \sigma_i \sigma_{i+2} \rangle = \frac{1}{4} \left[ (+1)(-1) + (+1)(-1) + (-1)(+1) + (-1)(+1) \right]
  = \frac{1}{4} (-1 -1 -1 -1) = -1.
  \]
%
Portanto,
%
\begin{equation}
    \varepsilon_{P4} = -J(0) + pJ(-1) = -pJ.
    \label{eq:eps_P4}
\end{equation}
%
    \item \textbf{Fase de período 3 (P3),} uma configuração com período 3, como \((+,+,-)\) repetida. Calculamos:
%
Sequência: \(+, +, -, +, +, -, \dots\)
%
Primeiros vizinhos: pares \((+,+)\), \((+,-)\), \((-,+)\), novamente. Média:
  \[
  \langle \sigma_i \sigma_{i+1} \rangle = \frac{1}{3} \left[ (+1)(+1) + (+1)(-1) + (-1)(+1) \right]
  = \frac{1}{3} (1 - 1 - 1) = -\frac{1}{3}.
  \]
%
Segundos vizinhos: pares \((+, -)\), \((+ , +)\), \((- , +)\). Média:
  \[
  \langle \sigma_i \sigma_{i+2} \rangle = \frac{1}{3} \left[ (+1)(-1) + (+1)(+1) + (-1)(+1) \right]
  = \frac{1}{3} (-1 + 1 - 1) = -\frac{1}{3}.
  \]
%
Portanto,
%
\begin{equation}
    \varepsilon_{P3} = -J\left(-\frac{1}{3}\right) + pJ\left(-\frac{1}{3}\right)
    = J\left( \frac{1}{3} - \frac{p}{3} \right)
    = \frac{J}{3}(1 - p).
    \label{eq:eps_P3}
\end{equation}
%
Temos as seguintes expressões para a energia por sítio (em unidades de \(J\)):
%
\[
\frac{\varepsilon_F}{J} = -1 + p,
\qquad
\frac{\varepsilon_{AF}}{J} = 1 + p,
\qquad
\frac{\varepsilon_{P4}}{J} = -p,
\qquad
\frac{\varepsilon_{P3}}{J} = \frac{1 - p}{3}.
\]
%
Vamos comparar cada par para determinar qual fase é a de menor energia para cada valor de \(p > 0\).
%
\paragraph{Comparação \(F\) vs \(P4\):}
%
\[
\varepsilon_F - \varepsilon_{P4} = J(-1 + p) - J(-p) = J(-1 + 2p).
\]
%
Portanto,
%
\[
\varepsilon_F < \varepsilon_{P4} \iff -1 + 2p < 0 \iff p < \frac{1}{2}.
\]
%
Para \(p < 1/2\), a fase ferromagnética é mais estável que a fase P4. Para \(p > 1/2\), a fase P4 é mais estável que a ferromagnética.
%
\paragraph{Comparação \(P4\) vs \(AF\):}
%
\[
\varepsilon_{P4} - \varepsilon_{AF} = J(-p) - J(1+p) = J(-1 - 2p) < 0 \quad \forall p > 0.
\]
%
Portanto, a fase P4 é sempre mais estável que a fase antiferromagnética.
%
\paragraph{Comparação \(P4\) vs \(P3\):}
\[
\varepsilon_{P4} - \varepsilon_{P3} = J(-p) - \frac{J}{3}(1-p)
= J\left( -p - \frac{1-p}{3} \right)
= J\left( \frac{-3p - 1 + p}{3} \right)
= J\left( \frac{-1 - 2p}{3} \right) < 0.
\]
%
Portanto, a fase P4 é sempre mais estável que a fase P3.
%
\paragraph{Comparação \(F\) vs \(AF\):}
%
\[
\varepsilon_F - \varepsilon_{AF} = J(-1+p) - J(1+p) = -2J < 0,
\]
%
logo a fase ferromagnética é sempre mais estável que a antiferromagnética.
%
\paragraph{Comparação \(F\) vs \(P3\):}
\[
\varepsilon_F - \varepsilon_{P3} = J(-1+p) - \frac{J}{3}(1-p)
= J\left( -1 + p - \frac{1}{3} + \frac{p}{3} \right)
= J\left( -\frac{4}{3} + \frac{4p}{3} \right)
= \frac{4J}{3}(p - 1).
\]
%
Portanto,
%
\[
\varepsilon_F < \varepsilon_{P3} \iff p < 1.
\]
%
Para \(p < 1\), a fase ferromagnética é mais estável que a P3; para \(p > 1\), a fase P3 é mais estável.
%
\paragraph{Comparação \(AF\) vs \(P3\):}
%
\[
\varepsilon_{AF} - \varepsilon_{P3} = J(1+p) - \frac{J}{3}(1-p)
= J\left( 1 + p - \frac{1}{3} + \frac{p}{3} \right)
= J\left( \frac{2}{3} + \frac{4p}{3} \right) > 0,
\]
%
logo a fase P3 é sempre mais estável que a AF.
%
Com base nas comparações acima, o estado fundamental é determinado da seguinte forma:
%
\begin{itemize}
    \item Para \(0 < p < 1/2\): a fase ferromagnética (F) é a de menor energia, pois \(\varepsilon_F < \varepsilon_{P4}\) e \(\varepsilon_F < \varepsilon_{P3}\).
    \item Para \(p = 1/2\): \(\varepsilon_F = \varepsilon_{P4} = -J/2\). Há uma degenerescência entre a fase ferromagnética e a fase de período 4. Este é o ponto de transição de primeira ordem entre as duas fases.
    \item Para \(1/2 < p < 1\): a fase de período 4 (P4) é a de menor energia, pois \(\varepsilon_{P4} < \varepsilon_F\) e \(\varepsilon_{P4} < \varepsilon_{P3}\).
    \item Para \(p = 1\): \(\varepsilon_{P4} = -J\) e \(\varepsilon_{P3} = 0\), então a fase P4 ainda é a de menor energia.
    \item Para \(p > 1\): a fase de período 3 (P3) torna-se a de menor energia, pois \(\varepsilon_{P3} < \varepsilon_{P4}\) (visto que para \(p > 1\), \(\varepsilon_{P4} - \varepsilon_{P3} = J(-1 - 2p)/3 < 0\)? Cuidado: a comparação P4 vs P3 mostrou que P4 é sempre menor. Vamos reavaliar: de fato, \(\varepsilon_{P4} - \varepsilon_{P3} = J(-1 - 2p)/3\), que é sempre negativo para \(p>0\). Portanto, P4 é sempre mais estável que P3! Isso contradiz a intuição de que para p grande a fase de período 3 poderia ser favorecida. Vamos verificar os cálculos.

    Refaçamos a comparação P4 vs P3:
    \[
    \varepsilon_{P4} = -pJ, \qquad \varepsilon_{P3} = \frac{J}{3}(1-p).
    \]
    Para \(p > 1\), temos \(\varepsilon_{P4} = -pJ\) e \(\varepsilon_{P3}\) é negativo, mas menos negativo? Por exemplo, para \(p = 2\): \(\varepsilon_{P4} = -2J\), \(\varepsilon_{P3} = -J/3\). Portanto, P4 é de fato menor (mais negativo). Então P4 sempre vence P3. Portanto, a fase P3 nunca é o estado fundamental para \(p>0\). Vamos verificar se há outras fases, como período 5, etc. Na verdade, para o modelo ANNNI unidimensional com interações de primeiro e segundo vizinhos, o estado fundamental é conhecido: para \(0 < p < 1/2\) é ferromagnético, para \(p > 1/2\) é a fase de período 4 (também chamada de fase "up-up-down-down" ou fase antiferro de segundo vizinho). Não há fase de período 3 no estado fundamental. Portanto, nossa análise está correta: P3 não compete.

    Vamos também verificar a fase de período 6, etc. Mas, em geral, a competição entre \(J_1 = -J\) (ferro) e \(J_2 = pJ\) (anti-ferro) leva a uma fase com período 4 para \(p > 1/2\). Isso é um resultado conhecido para a cadeia de Ising com interações de primeira e segunda vizinhança.
\end{itemize}
%
\end{itemize}
%
Portanto, o estado fundamental é:
%
\begin{equation}
    \boxed{
    \begin{cases}
        \text{Ferromagnético (F)}, & 0 < p < \frac{1}{2}, \\[4pt]
        \text{Degenerado entre F e P4}, & p = \frac{1}{2}, \\[4pt]
        \text{Período 4 (P4)}, & p > \frac{1}{2}.
    \end{cases}
    }
    \label{eq:gs_annni}
\end{equation}
%
A energia por sítio no estado fundamental é, portanto,
%
\begin{equation}
    \frac{\varepsilon_0}{J} =
    \begin{cases}
        -1 + p, & 0 < p \leq \frac{1}{2}, \\[4pt]
        -p, & p \geq \frac{1}{2}.
    \end{cases}
    \label{eq:eps0_annni}
\end{equation}
%
\begin{itemize}
    \item Para \(p < 1/2\), a interação ferromagnética de primeiros vizinhos domina, e o sistema se alinha ferromagneticamente.
    \item Para \(p > 1/2\), a interação antiferromagnética de segundos vizinhos torna-se suficientemente forte para frustrar o ferromagnetismo, e o sistema adota uma configuração de período 4, onde cada bloco de dois spins paralelos é seguido por dois spins antiparalelos.
    \item Em \(p = 1/2\), há uma competição exata entre as duas fases, resultando em uma degenerescência macroscópica do estado fundamental.
\end{itemize}
%
Este é um exemplo clássico de frustração geométrica em uma dimensão, onde a competição entre interações de diferentes alcances gera uma transição de fase de primeira ordem a \(T=0\).
%
No ponto \(p = 1/2\), ambas as expressões dão \(-1/2\), confirmando a degenerescência.
%
\subsection*{(ii) Matriz de transferência, autovalores e possibilidade de anomalia termodinâmica}
%
Como a interação alcança segundos vizinhos, a matriz de transferência usual (baseada em um único spin) não é suficiente; usamos como variável de estado o \emph{par} $\mu_i=(\sigma_i,\sigma_{i+1})$, com $4$ estados possíveis, e escrevemos
%
\begin{equation}
    T_{(\sigma,\sigma'),(\sigma',\sigma'')}=\exp\left[K\,\sigma'\sigma''-pK\,\sigma\sigma''\right],
\end{equation}
%
definida apenas quando o segundo índice do primeiro par coincide com o primeiro índice do segundo (condição de vizinhança). Na base ordenada $\{(+,+),(+,-),(-,+),(-,-)\}$, a matriz $4\times4$ resultante é
%
\begin{equation}
    T=
    \begin{pmatrix}
        \e{K(1-p)} & \e{K(p-1)} & 0 & 0\\
        0 & 0 & \e{-K(1+p)} & \e{K(1+p)}\\
        \e{K(1+p)} & \e{-K(1+p)} & 0 & 0\\
        0 & 0 & \e{K(p-1)} & \e{K(1-p)}
    \end{pmatrix}.
    \label{eq:T_annni}
\end{equation}
%
Essa matriz é invariante sob a simetria global $\sigma\to-\sigma$, que troca os estados $1\leftrightarrow4$ e $2\leftrightarrow3$. Nas combinações simétrica/antissimétrica $e_\pm=(\ket1\pm\ket4)/\sqrt2$, $f_\pm=(\ket2\pm\ket3)/\sqrt2$, a matriz $T$ bloco-diagonaliza em dois blocos $2\times2$ desacoplados,
%
\begin{align}
    M_A&=
    \begin{pmatrix}
        \e{K(1-p)} & \e{K(1+p)}\\
        \e{K(p-1)} & \e{-K(1+p)}
    \end{pmatrix}
    \quad\text{(setor }e_+,f_+\text{)},
    \\[4pt]
    M_B&=
    \begin{pmatrix}
        \e{K(1-p)} & -\e{K(1+p)}\\
        \e{K(p-1)} & -\e{-K(1+p)}
    \end{pmatrix}
    \quad\text{(setor }e_-,f_-\text{)}.
\end{align}
%
Os traços e determinantes são
%
\begin{align}
    \Tr M_A&=\e{K(1-p)}+\e{-K(1+p)}, & \det M_A&=-2\sinh(2Kp),\\
    \Tr M_B&=\e{K(1-p)}-\e{-K(1+p)}, & \det M_B&=+2\sinh(2Kp),
\end{align}
%
de modo que os quatro autovalores de $T$ são
%
\begin{equation}
    \lambda_{A,\pm}=\frac{\Tr M_A\pm\sqrt{(\Tr M_A)^2+8\sinh(2Kp)}}{2},
    \qquad
    \lambda_{B,\pm}=\frac{\Tr M_B\pm\sqrt{(\Tr M_B)^2-8\sinh(2Kp)}}{2}.
\end{equation}
%
Como $(\Tr M_A)^2+8\sinh(2Kp)>0$ sempre (para $p>0$), o setor $A$ tem autovalores reais garantidos. Pelo mesmo argumento de Perron-Frobenius do item anterior (a matriz~\eqref{eq:T_annni} tem entradas não negativas e é irredutível para $p>0$), o autovalor \emph{dominante} deve pertencer ao setor totalmente simétrico $e_+$:
%
\begin{equation}
    \boxed{
    \lambda_{\max}=\lambda_{A,+}
    =\frac12\left[\e{K(1-p)}+\e{-K(1+p)}+\sqrt{\left(\e{K(1-p)}+\e{-K(1+p)}\right)^2+8\sinh(2Kp)}\right].
    }
    \label{eq:lambda_max_annni}
\end{equation}
%
\noindent\textbf{Anomalia termodinâmica.} A energia livre por sítio é $f=-k_BT\ln\lambda_{\max}$, e a Eq.~\eqref{eq:lambda_max_annni} mostra que $\lambda_{\max}(K,p)$ é uma função \emph{analítica} de $K$ e $p$ para toda temperatura finita ($K<\infty$) -- não há nenhuma singularidade (zero do discriminante, cruzamento de níveis do autovalor dominante) em $T>0$. Isso é uma manifestação do teorema geral de ausência de transições de fase em $T>0$ para modelos unidimensionais com interações de alcance finito (argumento de Landau/van Hove): mesmo com a frustração introduzida pela competição $J$ vs.\ $pJ$, \emph{não há transição termodinâmica genuína em temperatura finita}.
%
No entanto, há sim um fenômeno interessante ligado à competição: o discriminante do setor $B$, $(\Tr M_B)^2-8\sinh(2Kp)$, pode tornar-se \emph{negativo} em certas regiões de $(p,K)$, e nesse caso $\lambda_{B,\pm}$ formam um par complexo conjugado. Embora isso não afete $\lambda_{\max}$ (que continua real e analítico), afeta o autovalor \emph{subdominante} que controla o decaimento a longas distâncias da função de correlação -- ver item (iii). A linha no plano $(p,t)$ onde esse discriminante se anula é a chamada \emph{linha de desordem} (\textit{disorder line}) do modelo ANNNI: não é uma transição de fase (a energia livre permanece analítica), mas marca a mudança qualitativa entre correlações que decaem monotonicamente e correlações que decaem de forma oscilatória/modulada -- o precursor, em $T>0$, da ordem antifase que só se torna genuína (com degenerescência macroscópica) exatamente em $T=0$, $p=1/2$, como visto no item (i).
%
\subsection*{(iii) Função de correlação de pares}
%
A correlação $\avg{\sigma_k\sigma_{k+r}}$ é obtida inserindo, na base de pares, o operador diagonal $\hat O=\mathrm{diag}(+1,+1,-1,-1)$ que mede o primeiro spin do par:
%
\begin{equation}
    \avg{\sigma_k\sigma_{k+r}}=\frac{\Tr\left[\hat O\,T^r\,\hat O\,T^{N-r}\right]}{\Tr\left[T^N\right]}
    \xrightarrow[N\to\infty]{}
    \sum_{n=2}^4 c_n\left(\frac{\lambda_n}{\lambda_{\max}}\right)^r,
    \label{eq:corr_annni_general}
\end{equation}
%
onde $\lambda_2,\lambda_3,\lambda_4$ são os autovalores subdominantes e $c_n$ são coeficientes (elementos de matriz de $\hat O$ entre autovetores à esquerda e à direita) fixados pela condição de normalização $\sum_n c_n=1$ (garantindo $\avg{\sigma_0\sigma_0}=1$). Diferentemente do problema (1), aqui $\hat O$ \emph{não} é diagonal na base que bloco-diagonaliza $T$ (ele troca $e_+\leftrightarrow e_-$ e $f_+\leftrightarrow f_-$), de modo que, em geral, todos os quatro autovalores contribuem para a soma em~\eqref{eq:corr_annni_general}. Ainda assim, o comportamento assintótico a longas distâncias é controlado apenas pelo autovalor de \emph{segunda maior norma}, $\lambda_2$:
%
\begin{equation}
    \avg{\sigma_0\sigma_r}\underset{r\gg1}{\sim} c_2\left(\frac{\lambda_2}{\lambda_{\max}}\right)^r.
\end{equation}
%
Há duas possibilidades, de acordo com a discussão do item (ii):
%
\begin{itemize}
    \item \textbf{Decaimento monotônico} (setor $B$ com discriminante positivo, tipicamente $p\lesssim\tfrac12$): $\lambda_2=\lambda_{A,-}$ ou $\lambda_{B,\pm}$ é real, e
    \begin{equation}
        \boxed{
        \avg{\sigma_0\sigma_r}\sim c_2\,\e{-r/\xi},
        \qquad
        \xi^{-1}(p,t)=\ln\left|\frac{\lambda_{\max}}{\lambda_2}\right|.
        }
    \end{equation}
    \item \textbf{Decaimento oscilatório} (setor $B$ com discriminante negativo, tipicamente $p\gtrsim\tfrac12$): $\lambda_{B,\pm}=|\lambda_2|\,\e{\pm iq}$ formam par complexo conjugado, e a combinação real correspondente dá
    \begin{equation}
        \boxed{
        \avg{\sigma_0\sigma_r}\sim A\left(\frac{|\lambda_2|}{\lambda_{\max}}\right)^r\cos(qr+\varphi),
        \qquad
        q=\arg(\lambda_2),
        }
    \end{equation}
    com uma modulação de período $2\pi/q$ que, para $p$ suficientemente acima de $1/2$, tende à modulação de período $4$ (antifase $\langle2\rangle$) característica do estado fundamental encontrado no item (i).
\end{itemize}
%
A linha $q=0\leftrightarrow q\neq0$ no plano $(p,t)$ é exatamente a linha de desordem mencionada no item (ii). A Fig.~\ref{fig:annni_corr} mostra $\avg{\sigma_0\sigma_r}$ calculada numericamente, a partir da diagonalização exata da matriz $4\times4$~\eqref{eq:T_annni}, para $p=0{,}2<1/2$ (decaimento puramente exponencial) e $p=1{,}5>1/2$ (decaimento oscilatório com período $4$, sinal claro da tendência à ordem antifase). O acordo com a discussão acima é excelente.
%
\begin{figure}[H]
    \centering
    \includegraphics[width=0.68\linewidth]{Figures/statistical_mechanics/annni_correlation.png}
    \caption{Função de correlação $\avg{\sigma_0\sigma_r}$ do modelo ANNNI, calculada por diagonalização exata da matriz de transferência~\eqref{eq:T_annni} para $K=\beta J=1{,}2$. Para $p<1/2$ o decaimento é puramente exponencial (monotônico); para $p>1/2$ surge uma modulação oscilatória de período $4$, precursora da ordem antifase do estado fundamental (item i).}
    \label{fig:annni_corr}
\end{figure}
% ============================================================
%
\section*{Problema (4) -- Teoria de Landau da transição ferromagnética}
%
A energia livre de Landau, na vizinhança imediata do ponto crítico, é
%
\begin{equation}
    f(T,m)=A+B(T)m^2+Cm^4+\mathcal O(m^6),
    \qquad
    B(T)=at=a\,\frac{T-T_c}{T_c},
    \qquad a,C>0.
    \label{eq:landau_f}
\end{equation}
%
onde $a$ e $C$ são constantes fenomenológicas, e $t = (T - T_c)/T_c$ é a temperatura reduzida. O termo de ordem $m^6$ é desprezado nas vizinhanças do ponto crítico, onde $m$ é pequeno.
%
\subsection*{(a) Esboço de $f(T,m)$ acima e abaixo de $T_c$}
%
Para $t > 0$, temos $B = a t > 0$. A energia livre $f-A=atm^2+Cm^4$ é uma função par, positiva para $m\neq0$, com um único mínimo em $m=0$ (poço único). A segunda derivada em $m = 0$ é
%
\begin{equation}
    \left. \frac{\partial^2 f}{\partial m^2} \right|_{m=0} = 2a t > 0,
\end{equation}
%
confirmando que $m = 0$ é um mínimo estável. Para $t < 0$, temos $B = a t = -a |t| < 0$. A energia livre é
%
\begin{equation}
    f - A = -a |t| m^2 + C m^4.
\end{equation}
%
Agora a função tem a forma de um poço duplo: $m = 0$ torna-se um máximo local, e dois mínimos simétricos aparecem em $m = \pm m_0 \neq 0$. A segunda derivada em $m = 0$ é
%
\begin{equation}
    \left. \frac{\partial^2 f}{\partial m^2} \right|_{m=0} = 2a t < 0,
\end{equation}
%
confirmando que $m = 0$ é um máximo local (instável). Os mínimos serão calculados no item (b).
%
A Fig.~\ref{fig:landau_f} ilustra esses dois regimes.
%
\begin{figure}[h]
\centering
\begin{tikzpicture}
\begin{axis}[
    width=0.45\textwidth,
    height=0.4\textwidth,
    xlabel={$m$},
    ylabel={$f - A$},
    xmin=-2, xmax=2,
    ymin=-0.5, ymax=4,
    grid=both,
    title={$t > 0$ (paramagnético)},
    legend pos=north east
]
\addplot[domain=-2:2, samples=100, color=blue, thick] {0.5*x^2 + 0.25*x^4};
\addlegendentry{$a t m^2 + C m^4$}
\end{axis}
\end{tikzpicture}
\hfill
\begin{tikzpicture}
\begin{axis}[
    width=0.45\textwidth,
    height=0.4\textwidth,
    xlabel={$m$},
    ylabel={$f - A$},
    xmin=-2, xmax=2,
    ymin=-1, ymax=1.5,
    grid=both,
    title={$t < 0$ (ferromagnético)},
    legend pos=north east
]
\addplot[domain=-2:2, samples=100, color=red, thick] {-0.5*x^2 + 0.25*x^4};
\addlegendentry{$-a|t| m^2 + C m^4$}
\end{axis}
\end{tikzpicture}
\caption{Energia livre de Landau para $t > 0$ (poço único, mínimo em $m=0$) e $t < 0$ (poço duplo, mínimos em $m = \pm m_0$).}
\label{fig:landau_f}
\end{figure}
%
\subsection*{(b) Expoente crítico $\beta$}
%
Minimizamos a energia livre em relação a $m$:
%
\begin{equation}
    \frac{\partial f}{\partial m} = 2a t m + 4 C m^3 = 2m (a t + 2 C m^2) = 0.
\end{equation}
%
As soluções são:
%
\begin{equation}
    m = 0, \qquad \text{ou} \qquad m^2 = -\frac{a t}{2C}.
\end{equation}
%
A segunda solução só é real para $t < 0$. Portanto, a magnetização espontânea de equilíbrio é
%
\begin{equation}
    m_0(t) = \pm \sqrt{-\frac{a t}{2C}} = \pm \sqrt{\frac{a}{2C}} \, |t|^{1/2}, \qquad t < 0.
\end{equation}
%
Comparando com a forma assintótica $m_0 \sim (-t)^\beta$, identificamos
%
\begin{equation}
    \boxed{\beta = \frac{1}{2}, \qquad m_0 = \pm \sqrt{\frac{a}{2C}} \, |t|^{1/2}}.
\end{equation}
%
A segunda derivada em $m = m_0$ é
%
\begin{equation}
    \frac{\partial^2 f}{\partial m^2} = 2a t + 12 C m^2 = 2a t + 12 C \left( -\frac{a t}{2C} \right) = 2a t - 6a t = -4a t = 4a |t| > 0,
\end{equation}
%
confirmando que $m = \pm m_0$ são mínimos estáveis e $m=0$ é um máximo local para $t<0$. A segunda derivada, vale $2at$ em $m=0$ (positiva, mínimo estável, para $t>0$. Esse é o resultado clássico de campo médio para o expoente do parâmetro de ordem -- idêntico ao valor $\beta=1/2$ obtido para o fluido de van der Waals na primeira lista, refletindo a universalidade da classe de campo médio.
%
\subsection*{(c) Calor específico a campo nulo}
%
A energia livre singular (a parte que depende de $m$) é obtida substituindo $m_0$ em $f$, substituindo $m_0$ de volta em $f$ (a menos da parte regular $A$, que tratamos como constante para isolar a contribuição \emph{singular} devida ao parâmetro de ordem, em analogia direta com a função $f_0(T)$ do problema 8 da primeira lista):
%
\begin{itemize}
    \item Para $t > 0$: $m_0 = 0 \Rightarrow f_{\rm sing} = 0$.
    \item Para $t < 0$:
    \begin{align}
        f_{\rm sing} &= a t m_0^2 + C m_0^4 \\
        &= a t \left( -\frac{a t}{2C} \right) + C \left( -\frac{a t}{2C} \right)^2 \\
        &= -\frac{a^2 t^2}{2C} + \frac{a^2 t^2}{4C} \\
        &= -\frac{a^2 t^2}{4C}.
    \end{align}
\end{itemize}
%
A entropia singular é
%
\begin{equation}
    s_{\rm sing} = -\frac{\partial f_{\rm sing}}{\partial T} = -\frac{1}{T_c} \frac{\partial f_{\rm sing}}{\partial t}.
\end{equation}
%
Para $t < 0$:
%
\begin{equation}
    s_{\rm sing} = -\frac{1}{T_c} \left( -\frac{a^2}{2C} t \right) = \frac{a^2 t}{2C T_c}.
\end{equation}
%
O calor específico singular é
%
\begin{equation}
    c_{\rm sing} = T \frac{\partial s_{\rm sing}}{\partial T} = T \frac{\partial s_{\rm sing}}{\partial t} \frac{\partial t}{\partial T}.
\end{equation}
%
Como $\partial t/\partial T = 1/T_c$, e próximo a $T_c$ temos $T \approx T_c$:
%
\begin{equation}
    c_{\rm sing} = \frac{T_c}{T_c} \frac{\partial s_{\rm sing}}{\partial t} = \frac{a^2}{2C T_c}, \qquad t < 0.
\end{equation}
%
Para $t > 0$, $f_{\rm sing} = 0$, logo $s_{\rm sing} = 0$ e $c_{\rm sing} = 0$. Portanto,
%
\begin{equation}
    \boxed{
    c_{\rm sing}(t) =
    \begin{cases}
        0, & t > 0, \\[4pt]
        \dfrac{a^2}{2C T_c}, & t < 0.
    \end{cases}
    }
\end{equation}
%
A descontinuidade no calor específico é
%
\begin{equation}
    \Delta c = \lim_{t \to 0^-} c_{\rm sing} - \lim_{t \to 0^+} c_{\rm sing} = \frac{a^2}{2C T_c}.
\end{equation}
%
Comparando com a forma $c \sim |t|^{-\alpha}$, identificamos
%
\begin{equation}
    \boxed{\alpha = 0}.
\end{equation}
%
\begin{figure}[h]
\centering
\begin{tikzpicture}
\begin{axis}[
    width=0.7\textwidth,
    height=0.5\textwidth,
    xlabel={$t$},
    ylabel={$c_{\rm sing}(t)$},
    xmin=-2, xmax=2,
    ymin=0, ymax=1.5,
    grid=both,
    title={Calor específico singular a campo nulo},
    legend pos=north east
]
\addplot[domain=-2:0, samples=100, color=red, thick] {1};
\addplot[domain=0:2, samples=100, color=blue, thick] {0};
\addlegendentry{$t < 0$}
\addlegendentry{$t > 0$}
\end{axis}
\end{tikzpicture}
\caption{Calor específico singular na teoria de Landau. Há uma descontinuidade finita em $t = 0$, sem divergência ($\alpha = 0$).}
\label{fig:landau_cv}
\end{figure}
%
$\alpha=0$ significa precisamente esse comportamento: sem divergência, mas com possível descontinuidade finita. A Fig.~\ref{fig:landau_cv} esboça esse comportamento -- o mesmo tipo de descontinuidade finita (sem divergência) encontrado para o fluido de van der Waals na primeira lista.
%
\subsection*{(d) Energia livre com campo externo e construção de Maxwell}
%
Na presença de um campo externo $H$ conjugado a $m$, a energia livre de Gibbs (ou a energia livre de Landau apropriada) é dada pela transformada de Legendre
%
\begin{equation}
    g(t, H) = \min_{m} \, g(t, H, m), \qquad g(t, H, m) = f(t, m) - H m.
\end{equation}
%
A condição de mínimo $\partial g/\partial m = 0$ fornece a equação de estado de Landau:
%
\begin{equation}
    \frac{\partial f}{\partial m} = H \quad \Longrightarrow \quad H = 2a t m + 4 C m^3.
    \label{eq:landau_eos}
\end{equation}
%
Esta equação é o análogo direto da equação de van der Waals para fluidos. A segunda derivada de $g$ em relação a $m$ é
%
\begin{equation}
    \frac{\partial^2 g}{\partial m^2} = \frac{\partial H}{\partial m} = 2a t + 12 C m^2.
\end{equation}
%
Para $t > 0$, $\partial^2 g/\partial m^2 > 0$ para todo $m$, logo $g$ é convexa e a solução $m(H)$ é única.
%
Para $t < 0$, a segunda derivada se anula em
%
\begin{equation}
    m_s^2 = -\frac{a t}{6C} = \frac{a |t|}{6C}.
\end{equation}
%
Entre $-m_s$ e $+m_s$, $\partial^2 g/\partial m^2 < 0$, indicando instabilidade termodinâmica. Esta região é substituída pela construção de Maxwell, que determina o campo de coexistência $H_{\rm coex}$. A isoterma $H(m)$ deixa de ser monotônica entre $-m_s$ e $+m_s$, exatamente como as ``alças'' de van der Waals no plano $p$--$v$.
%
A condição de áreas iguais (ou igualdade de pressões) é
%
\begin{equation}
    \int_{m_1}^{m_2} H(m) \, dm = H_{\rm coex} (m_2 - m_1),
\end{equation}
%
onde $m_1$ e $m_2$ são as magnetizações das duas fases coexistentes. Devido à simetria $m \to -m$, temos $m_1 = -m_0$ e $m_2 = +m_0$. Além disso, $H(m)$ é uma função ímpar:
%
\begin{equation}
    H(-m) = -H(m).
\end{equation}
%
Portanto,
%
\begin{equation}
    \int_{-m_0}^{m_0} H(m) \, dm = 0,
\end{equation}
%
pois a integral de uma função ímpar em um intervalo simétrico é zero. Logo,
%
\begin{equation}
    H_{\rm coex} (2m_0) = 0 \quad \Longrightarrow \quad H_{\rm coex} = 0.
\end{equation}
%
Assim, a coexistência entre as fases com magnetizações $-m_0$ e $+m_0$ ocorre em $H = 0$, para qualquer $t < 0$. A magnetização salta descontinuamente de $-m_0$ para $+m_0$ ao atravessar $H = 0$, caracterizando uma transição de primeira ordem induzida por campo.
%
\begin{figure}[h]
\centering
\begin{tikzpicture}
\begin{axis}[
    width=0.7\textwidth,
    height=0.5\textwidth,
    xlabel={$m$},
    ylabel={$H$},
    xmin=-2, xmax=2,
    ymin=-3, ymax=3,
    grid=both,
    title={Equação de estado de Landau para $t < 0$},
    legend pos=north west
]
\addplot[domain=-2:2, samples=200, color=blue, thick] {-x + x^3};
\addplot[domain=-1.2:1.2, samples=100, color=red, thick, dashed] {0};
\addplot[mark=*, mark size=3, color=black] coordinates {(-1,0) (1,0)};
\addlegendentry{$H = 2a t m + 4C m^3$}
\addlegendentry{$H = 0$ (coexistência)}
\end{axis}
\end{tikzpicture}
\caption{Isoterma $H(m)$ para $t < 0$. A região entre os extremos locais (tracejada) é instável; a construção de Maxwell, aqui trivial por simetria, fixa a coexistência em $H = 0$.}
\label{fig:landau_Hm}
\end{figure}
%
A construção de Maxwell (áreas iguais / dupla tangente, como no problema 3 da primeira lista) determina em que valor de $H$ ocorre a coexistência entre as duas fases $m<0$ e $m>0$. Aqui, porém, a simetria $m\to-m$, $H\to-H$ do problema \emph{fixa exatamente} essa construção
sem necessidade de resolver numericamente a condição de áreas iguais -- diferentemente do fluido de van der Waals, aqui a simetria do problema já garante $H_{\rm coex}=0$ para $t<0$. A Fig.~\ref{fig:landau_Hm} ilustra a isoterma $H(m)$ para $t<0$.
%
\subsection*{(e) Susceptibilidade a campo nulo}
%
A suscetibilidade magnética é definida como
%
\begin{equation}
    \chi = \left( \frac{\partial m}{\partial H} \right)_T.
\end{equation}
%
Diferenciando a equação de estado \eqref{eq:landau_eos} em relação a $H$ mantendo $t$ fixo:
%
\begin{equation}
    1 = (2a t + 12 C m^2) \frac{\partial m}{\partial H} \quad \Longrightarrow \quad \chi = \frac{1}{2a t + 12 C m^2}.
\end{equation}
%
\textbf{Para $t>0$} (regime paramagnético), $H=0\Rightarrow m=0$, logo
%
\begin{equation}
    \boxed{
    \chi(T,0)=\frac{1}{2at}\equiv C_+\,t^{-\gamma},
    \qquad
    \gamma=1,\quad C_+=\frac{1}{2a}.
    }
\end{equation}
%
\textbf{Para $t<0$} (regime ferromagnético), $H=0\Rightarrow m=m_0$, com $m_0^2=-B/2C=a|t|/2C$ (item b) e $B=at=-a|t|$, de modo que
%
\begin{equation}
    2B+12Cm_0^2=-2a|t|+12C\cdot\frac{a|t|}{2C}=-2a|t|+6a|t|=4a|t|,
\end{equation}
%
\begin{equation}
    \boxed{
    \chi(T,0)=\frac{1}{4a|t|}\equiv C_-\,(-t)^{-\gamma'},
    \qquad
    \gamma'=1,\quad C_-=\frac{1}{4a}.
    }
\end{equation}
%
A razão entre as amplitudes da suscetibilidade é
%
\begin{equation}
    \frac{C_+}{C_-} = \frac{1/(2a)}{1/(4a)} = 2.
\end{equation}
%
\begin{equation}
    \boxed{\frac{C_+}{C_-} = 2}.
\end{equation}
%
Esta é uma previsão universal da teoria de Landau (campo médio) para a classe de universalidade com parâmetro de ordem escalar.
%
Os dois expoentes coincidem, $\gamma=\gamma'=1$ (novamente o valor de campo médio, igual ao do fluido de van der Waals), mas as \emph{amplitudes} são diferentes: $C_+/C_-=2$. Essa razão universal de amplitudes ($=2$ em campo médio) é uma predição testável e característica da teoria de Landau.
%
\subsection*{(f) Relação de Escala de Rushbrooke $\alpha+2\beta+\gamma=2$}
%
Reunindo os resultados dos itens anteriores,
%
\begin{equation}
    \alpha=0,\qquad \beta=\frac12,\qquad\gamma=1
    \;\;\Longrightarrow\;\;
    \alpha+2\beta+\gamma=0+2\left(\frac12\right)+1=2,
\end{equation}
%
verificando exatamente a relação de escala de Rushbrooke,
%
\begin{equation}
    \boxed{\alpha+2\beta+\gamma=2.}
\end{equation}
%
Essa relação é uma consequência geral da hipótese de escala da energia livre singular (parte (g), abaixo) e vale, de forma universal, dentro de \emph{qualquer} classe de universalidade -- inclusive para fluidos simples, como o fluido de van der Waals da primeira lista, cujos expoentes $\alpha=0$, $\beta=1/2$, $\gamma=1$ também a satisfazem trivialmente, já que pertencem à mesma classe de universalidade de campo médio. A base microscópica dessa universalidade é o fato de que, próximo ao ponto crítico, o comprimento de correlação diverge e o sistema ``esquece'' os detalhes microscópicos (rede, forma exata das interações), retendo apenas a dimensionalidade, o número de componentes do parâmetro de ordem e o alcance das interações -- exatamente os ingredientes que definem uma classe de universalidade.
%
\subsection*{(g) Forma de escala da energia livre}
%
Buscamos a hipótese de escala (Widom) para a parte singular de $g(t, H)$. Usando a solução exata de Landau, minimizamos
%
\begin{equation}
    g(t, H, m) = a t m^2 + C m^4 - H m.
\label{eq:g_def}    
\end{equation}
%
Introduzimos a mudança de variável
%
\begin{equation}
    m = |t|^{1/2} \mu.
\end{equation}
%
Substituindo:
%
\begin{align}
    g(t, H, m) &= a t \, |t| \mu^2 + C |t|^2 \mu^4 - H |t|^{1/2} \mu \\
    &= a \, \text{sgn}(t) \, |t|^2 \mu^2 + C |t|^2 \mu^4 - H |t|^{1/2} \mu \\
    &= |t|^2 \left[ a \, \text{sgn}(t) \mu^2 + C \mu^4 - \frac{H}{|t|^{3/2}} \mu \right].
\end{align}
%
Minimizando sobre $\mu$, o resultado só pode depender da razão
%
\begin{equation}
    x \equiv \frac{H}{|t|^{3/2}}.
\end{equation}
%
Portanto, a energia livre singular assume a forma de escala generalizada
%
\begin{equation}
    \boxed{
    g_{\rm sing}(t, H) = |t|^{2} \, \Phi\!\left( \frac{H}{|t|^{3/2}} \right),
    }
\end{equation}
%
onde $\Phi$ é uma função de escala universal (dentro da teoria de Landau).
%
Equivalentemente, sob a transformação $t \to \lambda t$ e $H \to \lambda^{3/2} H$, a energia livre singular se comporta como uma função homogênea generalizada de grau $2$. Para isso, partimos da definição \eqref{eq:g_def} e aplicamos a mudança de variáveis $m = \lambda^{1/2} \tilde{m}$:
%
\begin{align}
    g_{\rm sing}(\lambda t, \lambda^{3/2} H)
    &= \min_{m} \left[ a (\lambda t) m^2 + C m^4 - (\lambda^{3/2} H) m \right] \\
    &= \min_{\tilde{m}} \left[ a (\lambda t) (\lambda^{1/2} \tilde{m})^2 + C (\lambda^{1/2} \tilde{m})^4 - (\lambda^{3/2} H) (\lambda^{1/2} \tilde{m}) \right] \\
    &= \min_{\tilde{m}} \left[ a (\lambda t) (\lambda \tilde{m}^2) + C (\lambda^2 \tilde{m}^4) - (\lambda^{3/2} H) (\lambda^{1/2} \tilde{m}) \right] \\
    &= \min_{\tilde{m}} \left[ \lambda^2 a t \tilde{m}^2 + \lambda^2 C \tilde{m}^4 - \lambda^2 H \tilde{m} \right] \\
    &= \lambda^2 \min_{\tilde{m}} \left[ a t \tilde{m}^2 + C \tilde{m}^4 - H \tilde{m} \right] \\
    &= \lambda^2 g_{\rm sing}(t, H).
\end{align}
%
\begin{equation}
    \boxed{
    g_{\rm sing}(\lambda t, \lambda^{3/2} H) = \lambda^{2} g_{\rm sing}(t, H),
    \qquad \text{para todo } \lambda > 0.
    }
\end{equation}
%
Na notação do enunciado, $g(\lambda^a t, \lambda^b H) = \lambda^{2-\alpha} g(t, H)$ com $\alpha = 0$, temos
%
\begin{equation}
    a = 1, \qquad b = \frac{3}{2}.
\end{equation}
%
Reconhecemos $b = \beta \delta$, onde $\delta$ é o expoente da isoterma crítica. Da equação de estado em $t = 0$:
%
\begin{equation}
    H = 4 C m^3 \quad \Longrightarrow \quad H \sim m^3 \quad \Longrightarrow \quad \delta = 3.
\end{equation}
%
Assim, $\beta \delta = (1/2) \cdot 3 = 3/2$, confirmando a consistência.
%
Esta forma de escala é a base microscópica (ou fenomenológica) por trás de todas as relações de escala entre expoentes críticos (Rushbrooke, Widom $\gamma = \beta(\delta - 1)$, etc.), pois basta diferenciar $g_{\rm sing}(t, H)$ o número apropriado de vezes em relação a $t$ e $H$ para obter $c_V$, $m$ e $\chi$, e comparar os expoentes resultantes.
%
\begin{table}[H]
\centering
\begin{tabular}{@{}lccc@{}}
\toprule
Grandeza & Definição & Resultado & Valor \\
\midrule
Calor específico & $c\sim|t|^{-\alpha}$ & salto finito & $\alpha=0$ \\
Parâmetro de ordem & $m\sim(-t)^{\beta}$ & $t<0,\,H=0$ & $\beta=1/2$ \\
Susceptibilidade & $\chi\sim|t|^{-\gamma}$ & $H=0$ & $\gamma=1$ \\
Isoterma crítica & $H\sim m^{\delta}$ & $t=0$ & $\delta=3$ \\
\bottomrule
\end{tabular}
\caption{Expoentes críticos de campo médio (Landau) obtidos neste problema -- idênticos aos do fluido de van der Waals (primeira lista), como esperado pela universalidade de campo médio.}
\end{table}
%%%%%%%%%%%%%%%%%%%%%%%%%%%%%%%%%%%%%%%%%%%%%%%%%%%%%%%%%%%%%%%%%%%%%%%%%%%%%%%%
\begin{thebibliography}{99}
\bibitem{domingues2026statmech}
Thiago S. Domingues, \textit{Statistical Mechanics: Repository to add materials about a statistical mechanics course}, GitHub repository, 2026. 
\url{https://github.com/ThiagoSDomingues/Statistical_Mechanics}
\bibitem{salinas} S. R. A. Salinas, \emph{Introdução à Física Estatística}, 2ª ed., 3ª reimpressão, Edusp, São Paulo, 2013.
\bibitem{krinsky} S. Krinsky and D. Furman, \emph{Phys. Rev. Lett.} \textbf{32}, 731 (1974).
\bibitem{stephenson} J. Stephenson, \emph{Can. J. Phys.} \textbf{48}, 129 (1970).
\bibitem{selke} W. Selke, \emph{The ANNNI model -- theoretical analysis and experimental application}, Phys. Rep. \textbf{170}, 213 (1988).
\bibitem{Gantmacher1990}
F. R. Gantmacher, \textit{The Theory of Matrices}, Vol.~2, Chelsea Publishing Company, New York (1990).
\bibitem{HornJohnson2013}
R. A. Horn and C. R. Johnson, \textit{Matrix Analysis}, 2nd ed., Cambridge University Press, Cambridge (2013).
\bibitem{MacCluer2000}
C. R. MacCluer, SIAM Rev. \textbf{42}, 487 (2000).
\end{thebibliography}
%%%%%%%%%%%%%%%%%%%%%%%%%%%%%%%%%%%%%%%%%%%%%%%%%%%%%%%%%%
\end{document}
